%% --------------------------------------------------------------------
%% Template UTEC Tesis
%% --------------------------------------------------------------------
% Este es el archivo que se debe compilar.
% 
% Este template ha sido modificado y actualizado por Eduardo Castro y Roosevelt Ubaldo en base a lo trabajado por Víctor Murray, Oscar Ramos y Juan Carlos Barbaran.
%
% Última actualización: Mayo, 2022

\documentclass[a4paper,12pt,oneside]{tesisutec}

\selectlanguage{spanish}

%% Paquetes
\usepackage[utf8]{inputenc}
\usepackage[square, numbers, comma, sort&compress]{natbib}
% Libreria de idioma
\usepackage[spanish]{babel}
% Libreria para posicionamiento
\usepackage{float}
% Librerias para insertar codigos
\usepackage[spanish,onelanguage,ruled,vlined]{algorithm2e}
\usepackage{verbatim} 
% Librería para hipervínculos
\usepackage{hyperref}
 % Librería necesaria para arreglar el orden de referencias en overleaf.com
\usepackage{notoccite}

% Incluir acá los paquetes adicionales que deseas. 
% Ubicación de los imágenes.
\graphicspath{images/}

\makeatletter
% Reinsert missing \algbackskip
\def\algbackskip{\hskip-\ALG@thistlm}
\makeatother

\hypersetup{urlcolor=blue, colorlinks=true}

%% ----------------------------------------------------------------------------
%% Marcas de GAP de redaccion (ver overleaf/CLAUDE.md y redaccion/ESTILO.md).
%% Hacen visible en el PDF lo que el texto todavia no puede afirmar.
%%   \GAPLIT{...}   falta una fuente: busqueda de literatura pendiente.
%%   \GAPDATO{...}  falta un resultado propio: experimento o revision pendiente.
%%   \GAPDEC{...}   falta una decision de la autora o del asesor.
%% Para la entrega, cambiar a \entregafinaltrue: cualquier GAP detiene la
%% compilacion con un error que dice cual es.
%% ----------------------------------------------------------------------------
\usepackage{xcolor}
\newif\ifentregafinal
\entregafinalfalse
\newcounter{gaplit}
\newcounter{gapdato}
\newcounter{gapdec}
\newcommand{\gapmarca}[4]{%
  \stepcounter{#1}%
  \ifentregafinal\errmessage{#2 sin resolver: #4}\fi
  \textcolor{#3}{\textbf{[#2: #4]}}}
\newcommand{\GAPLIT}[1]{\gapmarca{gaplit}{GAP-LIT}{red!75!black}{#1}}
\newcommand{\GAPDATO}[1]{\gapmarca{gapdato}{GAP-DATO}{orange!85!black}{#1}}
\newcommand{\GAPDEC}[1]{\gapmarca{gapdec}{GAP-DEC}{blue!70!black}{#1}}
\makeatletter
\AtEndDocument{\immediate\write\@auxout{%
  \string\gdef\string\gapstotales{\arabic{gaplit}/\arabic{gapdato}/\arabic{gapdec}}}}
\newcommand{\resumengaps}{%
  \ifentregafinal\else
    \@ifundefined{gapstotales}{}{%
      \begin{center}\fbox{\parbox{0.85\linewidth}{\centering\small
        \textbf{BORRADOR} --- GAP pendientes (literatura / dato / decisión):
        \textbf{\gapstotales}}}\end{center}}%
  \fi}
\makeatother

\begin{document}

\frontmatter

\department{CIENCIAS DE LA COMPUTACIÓN}
\degree{Bachiller }
\major{Ciencias de la Computación }


\title{Síntesis de imágenes de TC pélvica con implantes metálicos mediante difusión condicionada para aumento de datos en segmentación ósea}

\author{Kiara Alexandra Balcázar Santa Cruz \orcid{0000-0000-0000-0000}} % Es obligatorio agregar ORCID del alumno
\supervisor{Victor \orcid{0000-0000-0000-0000}} % Es obligatorio agregar ORCID del asesor
\date{2026}

\maketitle

\setstretch{1.5}

\resumengaps


%% ============================================================================
%  Dedicatoria
%% ============================================================================

\dedicatory{
xxxxxxxxxxxxxxxxxxxxxxxxxxxxxxxxxxxxxxxxxxxxxxxxxxx
xxxxxxxxxxxxxxxxxxxxxxxxxxxxxxxxxxxxxxxxxxxxxxxxxxx
xxxxxxxxxxxxxxxxxxxxxxxxxxxxxxxxxxxxxxxxxxxxxxxxxxx
xxxxxxxxxxxxxxxxxxxxxxxxxxxxxxxxxxxxxxxxxxxxxxxxxxx
xxxxxxxxxxxxxxxxxxxxxxxxxxxxxxxxxxxxxxxxxxxxxxxxxxx
XXxxxxxxxxxxxxxxxxxxxxxxxxxxxxxxxxxxxxxxxxxxxxxxxx
}

%% ============================================================================
%  Pagina de Agradecimientos
%% ============================================================================

\acknowledgements{

xxxxxxxxxxxxxxxxxxxxxxxxxxxxxxxxxxxxxxxxxxxxxxxxxxx
xxxxxxxxxxxxxxxxxxxxxxxxxxxxxxxxxxxxxxxxxxxxxxxxxxx
xxxxxxxxxxxxxxxxxxxxxxxxxxxxxxxxxxxxxxxxxxxxxxxxxxx
xxxxxxxxxxxxxxxxxxxxxxxxxxxxxxxxxxxxxxxxxxxxxxxxxxx
xxxxxxxxxxxxxxxxxxxxxxxxxxxxxxxxxxxxxxxxxxxxxxxxxxx
xxxxxxxxxxxxxxxxxxxxxxxxxxxxxxxxxxxxxxxxxxxxxxxxxxx

}



\tableofcontents

\newpage

\listoftables

\newpage

\listoffigures

\addtocontents{toc}{\vspace{1.5em}}

%% ============================================================================
\mainmatter
\pagestyle{fancy}

\customchapter{RESUMEN}

Un resumen debe ser un texto informativo, breve, claro, objetivo y veraz y debe permitir a los lectores entender el contenido del documento rápidamente. Un resumen es como una versión en miniatura de su trabajo y debe entenderse sin leer el documento. No contiene figuras, tablas, citas, referencias, abreviaturas y siglas. Se recomienda usar la estructura IMRyC:
\begin{itemize}
  \item \textbf{Introducción a la Investigación:} Sección que sintetiza qué es lo que se conocía sobre el problema de investigación, antecedentes y el propósito de la investigación. ¿Qué problema se intentó resolver? ¿Cuál era el objetivo de la investigación? ¿Por qué era importante este problema y los resultados de la investigación? \textit{(2-3 oraciones)}.
  \item \textbf{Marco Metodológico:} Usualmente incluye caracterización de la investigación y la declaración del cómo se resolvió el problema. ¿Cómo procedió en resolver el problema?. No se describen métodos, ni instrumentos ni variables, pero la declaración debe ser suficientemente precisa para transmitir el diseño y el perfil de la investigación \textit{(3-4 oraciones)}.
  \item \textbf{Resultados:} El texto debe ser descriptivo y suficientemente informativo de los resultados más significativos en el logro de los objetivos. ¿Cuál fue la respuesta? Exponga aquí los resultados en números y cantidades. Incluya en el resumen los hallazgos más importantes \textit{(4-6 oraciones)}.
  \item \textbf{Conclusiones:} A través del mensaje breve y preciso, se transmite la interpretación final de los logros de la investigación y otros hallazgos importantes. ¿Cuáles son las implicaciones de sus respuestas? Una correcta redacción de esta parte motivará al lector a seguir leyendo \textit{(2-3 oraciones)}.
\end{itemize}



\noindent \textbf{Palabras clave:}\\
\noindent Palabra clave 1; Palabra clave 2; Palabra clave 3; Palabra clave 4
\customchapter{ABSTRACT} 

\begin{center}
\large \vspace{-1.5cm} \textbf{TÍTULO DE LA TESIS EN INGLÉS}
\end{center}

Lorem ipsum dolor sit amet, cum tollit aeterno ut. Sea amet iuvaret placerat in, saperet impedit consulatu mei ne, quo no alii vituperatoribus. No mutat dolorum habemus pro, vis rebum accusam no, meis partiendo necessitatibus cu sit. Zril mediocrem scriptorem cu quo, ex duo mucius doming legendos. Posse nemore cu cum. Et facer pertinax sed, atqui ignota antiopam in pro, dolore oportere ius ex.

An justo nonumy mnesarchum quo, eruditi cotidieque in sed. Consulatu urbanitas conceptam eam id, liber accusamus eu mel. Ius iusto albucius elaboraret id, dico partiendo contentiones nec in, sint eruditi per ne. No eruditi probatus cum, no mei sumo facer torquatos. Sea dicunt facilisis corrumpit at, audiam accusam eligendi qui ne.

Facete officiis invenire sit ex, facete causae fabellas ei qui. Et pertinacia contentiones definitiones nam. Cum moderatius signiferumque in, blandit senserit cum no. Mea appareat posidonium contentiones an, no case accommodare cum, ei sea oporteat percipitur. 
Ea congue necessitatibus sed, no adhuc graeco usu. Eam harum equidem fierent ne, duo id facer consequuntur signiferumque. 
Dolore molestie phaedrum nec ei, ei eos noluisse constituto. Eruditi propriae in mea, vero laboramus cotidieque ad duo. Tacimates constituto est ex, ad mei clita sapientem evertitur. \\


\noindent \textbf{Keywords:}\\
\noindent Keyword 1; Keyword 2; Keyword 3; Keyword 4


 
\customchapter{INTRODUCCIÓN}

La segmentación automática de estructuras óseas en tomografías computarizadas (TC) pélvicas constituye una tarea relevante para el análisis de lesiones, la reconstrucción tridimensional y la planificación de procedimientos quirúrgicos. Los métodos basados en aprendizaje profundo han permitido automatizar esta tarea; sin embargo, su desempeño depende de que los datos de entrenamiento representen las variaciones que pueden encontrarse en las imágenes reales. Entre estas variaciones se encuentran las fracturas, las diferencias entre protocolos de adquisición y la presencia de implantes metálicos, siendo esta última una de las condiciones más difíciles para la segmentación de los huesos pélvicos \cite{liu2021ctpelvic1k}.

Los implantes metálicos, como tornillos, placas y prótesis, presentan una elevada atenuación de los rayos X y producen artefactos que suelen manifestarse como bandas oscuras, regiones brillantes y trazos en forma de rayas. Estos artefactos inducidos por metal pueden ocultar los límites de las estructuras anatómicas y reducir la calidad de las imágenes utilizadas por los modelos de segmentación \cite{selles2024marreview,xie2024implantsegmentation}. Su formación está relacionada con fenómenos físicos y computacionales, entre ellos el endurecimiento del haz, la escasez de fotones y las inconsistencias introducidas durante la reconstrucción. Por esta razón, reproducir de manera exacta el proceso de formación de los artefactos requiere normalmente información del dominio de proyección, como los sinogramas y los parámetros de adquisición.

La disponibilidad de datos anotados bajo esta condición es limitada. CTPelvic1K reúne 1\,184 volúmenes de TC pélvica procedentes de diferentes fuentes, de los cuales 75 contienen artefactos metálicos. El conjunto proporciona anotaciones óseas para 1\,109 volúmenes sin metal y para únicamente 14 volúmenes afectados por metal, mientras que los 61 casos restantes con metal se mantienen sin anotaciones óseas y no se utilizan en el entrenamiento supervisado original \cite{liu2021ctpelvic1k}. Esta distribución produce una representación reducida de los casos metálicos durante el entrenamiento y limita el desarrollo de modelos robustos ante este tipo de alteración.

El aumento de datos mediante transformaciones geométricas o fotométricas convencionales permite incrementar la variabilidad de un conjunto de entrenamiento, pero no modela explícitamente la presencia, la forma ni la ubicación de los implantes metálicos. Tampoco reproduce los patrones de intensidad asociados con los artefactos que estos producen. En este contexto, los modelos de difusión ofrecen un mecanismo para aprender distribuciones complejas y sintetizar imágenes médicas \cite{kazerouni2023diffusionsurvey}. Los modelos de difusión latente reducen el costo computacional al realizar el proceso de difusión en un espacio latente \cite{rombach2022latentdiffusion}, mientras que ControlNet permite incorporar condiciones espaciales, como mapas de segmentación o máscaras, en un modelo de difusión previamente entrenado \cite{zhang2023controlnet}.

La síntesis condicionada también permite evaluar la utilidad de las imágenes generadas mediante una tarea posterior (\textit{downstream task}). En síntesis de tumores, por ejemplo, se ha evaluado si los casos artificiales contribuyen al entrenamiento de modelos capaces de detectar o segmentar lesiones reales \cite{chen2024tumorsynthesis}. Esta forma de evaluación resulta pertinente porque una imagen sintética puede presentar similitud visual sin contener características útiles para un modelo de análisis médico. Esa evaluación presupone, sin embargo, que lo sintetizado sea plausible, y este trabajo examina esa plausibilidad antes que la utilidad (véase la justificación).

A partir de esta necesidad, la presente investigación propone una cadena de síntesis de dos componentes que inserta un tornillo de osteosíntesis y su artefacto metálico local en volúmenes de TC pélvica sin implante. El primer componente, un muestreador de colocación, propone la pose tridimensional del tornillo dentro del corredor óseo medido en cada volumen. El segundo, un sintetizador por difusión, genera su apariencia en el dominio de imagen, sin espacio latente ni datos de proyección. La generación se restringe a la máscara del implante y a una banda a su alrededor, para que el artefacto pueda aparecer fuera del metal, y el resto del volumen se copia sin cambios.

La evaluación mide la coherencia física y quirúrgica de lo sintetizado, y no su utilidad para entrenar una red de segmentación. La coherencia de la colocación se cuantifica comparando la distribución de grados de brecha cortical de las poses, es decir, de cuánto sobresale el implante del hueso, con distribuciones clínicas publicadas. La coherencia de la apariencia se cuantifica con las métricas de un protocolo publicado de simulación física, frente a ese protocolo y frente a una inserción por copia y pegado, que coloca el metal sin artefacto. Antes del sintetizador, una compuerta examina en qué representación de los valores de TC, expresados en unidades Hounsfield (HU), puede operar la síntesis.

El documento presenta, después de este capítulo, los fundamentos teóricos y el estado del arte relacionados con la tomografía, los artefactos inducidos por metal, la fijación iliosacra y los modelos de difusión condicionada. Posteriormente, se describe la metodología propuesta, el diseño experimental y las métricas de evaluación. Finalmente, se presentan y discuten los resultados, seguidos de las conclusiones, limitaciones y posibles líneas de trabajo futuro.

\introsection{Formulación del problema}

¿En qué medida la cadena propuesta genera colocaciones coherentes con las distribuciones clínicas de brecha cortical y artefactos coherentes con el protocolo físico de Peters et al.~\cite{peters2025hybrid}, frente a la inserción por copia y pegado?

\introsection{Objetivos de investigación}

% Fuentes de esta subseccion y de las dos siguientes: tesis/main.tex (Objectives, Out of scope,
% Recognized Gap, Datasets), docs/00-tesis.md, docs/01-decisiones.md y capitulo3.tex.
% Fuente de cada cifra: redaccion/rondas/introduccion-r00-respuesta.md e introduccion-r01-respuesta.md.
Los objetivos son los cuatro que desarrolla el Capítulo~\ref{cap:propuesta}, pero solo los Objetivos~2 y~3 son componentes de la cadena propuesta. El Objetivo~1 es una prueba previa a la cadena. Los Objetivos~1 a~3 tienen en ese capítulo una variable dependiente que los verifica; el Objetivo~4 no tiene variable dependiente propia, porque define o adopta las métricas con que se evalúan los Objetivos~2 y~3.

\paragraph{Objetivo general}

Diseñar, implementar y evaluar una cadena de síntesis por difusión en el dominio de imagen que genere implantes metálicos de osteosíntesis y su artefacto metálico local en volúmenes de TC pélvica, con coherencia física y quirúrgica. Los implantes de osteosíntesis son dispositivos de fijación ósea, como tornillos y placas. La colocación del implante se restringe con el corredor óseo medido sobre cada volumen, es decir, la región de hueso por la que puede pasar un tornillo sin atravesar la cortical. La apariencia se condiciona con una codificación multiventana de los HU, formada por varios canales, cada uno con una ventana de HU distinta. El único implante que se sintetiza es un tornillo transilíaco-transsacro, o transiliosacro (véase el alcance).

\paragraph{Objetivos específicos}

\begin{enumerate}
    \item \textbf{Evaluar si el autoencoder de un modelo de difusión latente conserva los HU de una codificación multiventana (compuerta \emph{Go/No-Go}).} Determinar si el paso de ida y vuelta (HU $\rightarrow$ representación $\rightarrow$ HU) por ese autoencoder conserva los HU del hueso con un error absoluto medio menor que 25~HU. El autoencoder es la red que comprime la imagen a un espacio latente, donde opera la difusión, y la reconstruye desde él \cite{rombach2022latentdiffusion}. La regla de la compuerta examina seis combinaciones, de dos variantes del autoencoder y tres codificaciones, sobre 34 pacientes de prueba, y aprueba si alguna cumple el criterio (Sección~\ref{sec:obj1}).

    \item \textbf{Diseñar un muestreador de colocación quirúrgicamente restringido.} Proponer poses tridimensionales de un tornillo transilíaco-transsacro, que cruza las dos articulaciones sacroilíacas de un ilion al otro y se modela como un cilindro paramétrico y rígido, alrededor del corredor óseo medido en cada volumen. Las poses se expresan en el marco de referencia de Kaiser et al.~\cite{kaiser2014dysmorphism}, y un corredor es viable si admite el calibre del tornillo más una holgura radial tomada de esa misma fuente. Cada pose recibe un \textbf{grado de brecha cortical}, es decir, su nivel en una escala ordinal de cuatro niveles según cuánto sobresale el implante del hueso. La distribución de grados se compara con la \textbf{referencia clínica}, formada por las series con navegación computarizada (navegada) y sin ella (convencional) de Zwingmann et al.~\cite{zwingmann2009navigated} en S1, el primer segmento del sacro. La comparación usa la distancia de Wasserstein-1, que sobre grados ordinales suma las diferencias absolutas entre las dos distribuciones acumuladas.

    \item \textbf{Diseñar un sintetizador condicionado.} Formular la síntesis como \emph{inpainting} por difusión 2.5D en el dominio de imagen, es decir, sobre los valores de TC por vóxel, sin espacio latente ni datos de proyección. El \emph{inpainting} regenera solo una región enmascarada de la imagen y conserva el resto; el término 2.5D se precisa en el marco teórico. Lo que se genera es la \textbf{región de generación} $G = M \cup B_{\delta}$, donde $M$ es la \textbf{máscara del implante} y $B_{\delta}$ una \textbf{banda de generación extendida} (\emph{generation band}) de unos 12~mm alrededor de ella. La banda permite que las rayas (\emph{streaking}) y las zonas oscuras del endurecimiento del haz aparezcan fuera del metal. La apariencia generada se compara con la \textbf{inserción por copia y pegado}, que coloca el metal sin artefacto, con el protocolo físico de Peters et al.~\cite{peters2025hybrid} y con observaciones reales de artefacto; en estas últimas, con estadísticos de los HU alrededor del metal que no exigen una imagen sin metal del mismo paciente (Sección~\ref{sec:apariencia}).

    \item \textbf{Formalizar la métrica de colocación y adoptar las métricas de apariencia de Peters et al.} Formalizar la admisibilidad quirúrgica de la colocación (SAP, \emph{Surgical Admissibility of Placement}), la única métrica que introduce este trabajo. SAP reúne dos componentes, definidos en el Objetivo~2: el grado de brecha cortical de cada pose y la viabilidad del corredor, que se evalúa con un calibre de 7.0~mm y una holgura radial de 1~mm (Sección~\ref{sec:sap}). Solo la distribución de grados se compara con la referencia clínica; la viabilidad se reporta de forma descriptiva. La definición de SAP incluye la distancia de Wasserstein-1 de esa comparación, que se ejecuta en el Objetivo~2. La apariencia se evalúa con las métricas de Peters et al., con sus nombres publicados: \emph{bone integrity}, \emph{metal integrity} y \emph{streak amplitude}. Esa adopción no se presenta como contribución evaluada. Esas métricas se diseñaron para evaluar la reducción de artefactos metálicos (MAR, \emph{metal artifact reduction}), y usarlas para evaluar síntesis exige invertirlas. La inversión está definida para la \emph{streak amplitude}; la \emph{bone integrity} y la \emph{metal integrity} conservan su fórmula e invierten su lectura, cuya referencia es el valor del protocolo físico (Sección~\ref{sec:apariencia}).
\end{enumerate}

La compuerta del Objetivo~1 se formuló para una síntesis por difusión latente guiada con ControlNet. Su veredicto fue negativo, y después de conocerlo el Objetivo~3 se reformuló en el dominio de imagen, sin modificar la regla de la compuerta. Con la misma regla y sobre los mismos pacientes de prueba se examinó después un autoencoder entrenado con TC (Sección~\ref{sec:obj1}). El veredicto descartó la ruta latente, no la codificación multiventana, que el sintetizador sigue usando, ahora sin la etapa del autoencoder. La ida y vuelta de esa codificación sin autoencoder se midió en los experimentos del Objetivo~1, aunque no entra en su criterio de fallo (Sección~\ref{sec:obj1}). Sobre los 34 pacientes de prueba, la media del error absoluto medio en el hueso fue de 0.00~HU en dos de las tres codificaciones y de 21.80~HU en la tercera. El sintetizador adopta una de esas dos, la de compresión arcoseno hiperbólico, en coma flotante de 32 bits (Sección~\ref{sec:sintetizador}).

No todos los objetivos fijan de antemano qué resultado los haría fallar; el Objetivo~1 sí lo fija. Su regla se escribió después de una exploración que incluía a los pacientes de prueba y cuyo error ya superaba el criterio, pero antes de la prueba que decide. El criterio de 25~HU es anterior a esa exploración y no se movió (Sección~\ref{sec:obj1}). El Objetivo~2 compara su distribución de grados con la referencia clínica sin regla de decisión, que fijada ahora sería posterior al resultado, y la lee contra la distancia entre las dos series de esa referencia (Sección~\ref{sec:sap}). En el Objetivo~3 decide la comparación de equivalencia frente al protocolo físico: solo se concluye equivalencia si el intervalo de confianza del 90\,\% de la diferencia pareada cae entero en un margen $[-\Delta, +\Delta]$, aún sin medir (Sección~\ref{sec:apariencia}). La superioridad frente a la inserción por copia y pegado, que no produce rayas por construcción, es un control de cordura que supera cualquier brazo que genere algo, y ningún resultado suyo cuenta a favor del sintetizador. El Objetivo~4 es un objetivo de definición: se verifica con los controles de implementación de SAP (identidad, monotonía y resolución), se ejerce en el Objetivo~2 y no se evalúa experimentalmente por separado.

\introsection{Justificación}

Este trabajo se justifica por una carencia de datos anotados. La segmentación ósea automática de la TC pélvica depende de volúmenes anotados, y los artefactos metálicos pueden ocultar los límites anatómicos cerca de los implantes. En CTPelvic1K, el conjunto que usa este trabajo, Liu et al.~\cite{liu2021ctpelvic1k} publican anotaciones óseas para 14 de los 75 volúmenes con artefacto metálico y dejan sin anotar los 61 restantes. La motivación de fondo es usar implantes sintéticos como aumento de datos para la segmentación ósea alrededor del implante, pero esa utilidad no se mide aquí (véase el alcance). Este trabajo evalúa, en su lugar, si lo sintetizado es coherente con la física del artefacto y con la práctica quirúrgica, porque medir la utilidad para segmentación de ejemplos cuya plausibilidad física no se conoce confundiría las dos propiedades.

La justificación computacional es que ninguno de los tres enfoques revisados (simulación física, planificación quirúrgica determinista y síntesis de lesiones por difusión) produce a la vez una distribución de poses del implante y la apariencia de su artefacto. Ninguno de los dos primeros produce esa distribución de poses. En la simulación física, el protocolo físico de Peters et al.~\cite{peters2025hybrid} simula la formación del artefacto, pero coloca el metal al azar en sus datos de entrenamiento, porque la colocación manual les resultaba impracticable. En su evaluación, sobre conjuntos de pacientes que el artículo llama clínicamente representativos, el metal se inserta en ubicaciones que llama pertinentes (\emph{meaningful locations}), sin que describa la regla con que se eligen. Los planificadores quirúrgicos deterministas, como el de Liu et al.~\cite{liu2025pipeline}, optimizan una única trayectoria. Zwingmann et al.~\cite{zwingmann2009navigated}, en cambio, reportan distribuciones de malposición distintas según la técnica quirúrgica, y este trabajo lee ese resultado como indicación de que el muestreador ha de reproducir una distribución de poses y no una pose única.

Este trabajo prevé ejecutar, además, el protocolo de Peters et al. sobre las poses del muestreador, en los pacientes de prueba sin metal y según el plazo, como comparación para la apariencia (Sección~\ref{sec:apariencia}). Ese protocolo no sustituye al sintetizador: necesita simular proyecciones para cada pose, simula en dos dimensiones una sola fila de detector y los objetos metálicos de sus datos de entrenamiento son formas fractales \cite{peters2025hybrid}. El sintetizador, en cambio, opera sobre cualquier TC reconstruida, sin proyecciones. Si alcanza la equivalencia con la simulación física es justamente lo que pone a prueba el Objetivo~3, de modo que esa equivalencia es una pregunta de este trabajo y no una premisa.

El tercer enfoque, la síntesis de lesiones por difusión, no liga al objeto generado los cambios de intensidad fuera de su máscara, y el artefacto metálico no queda dentro del implante. De Man et al.~\cite{deman1999} identifican el endurecimiento del haz, la dispersión, el ruido y el efecto exponencial de gradiente de borde como las causas principales de las rayas del metal, y las muestran irradiando desde el metal. Los modelos de difusión revisados generan la lesión por \emph{inpainting} acotado a su máscara \cite{ramzan2026claim,chen2024tumorsynthesis,jacob2026lgesynthnet} o sintetizan el corte entero desde ruido \cite{zhang2025diffboost}. En los que acotan el \emph{inpainting} a la máscara no hay un mecanismo para los cambios de intensidad que el objeto generado induce fuera de ella, y en ninguno hay una señal de entrenamiento para esos cambios.

La contribución de este trabajo es la combinación de tres elementos que las fuentes revisadas no reúnen. El primero son geometrías de implante rígidas y paramétricas como máscara de síntesis, que evitan segmentar con un umbral fijo de HU el metal que se coloca. Xie et al.~\cite{xie2024implantsegmentation} reportan que esa segmentación, en cortes simulados, cubrió el implante en exceso, y en la cohorte local los objetos metálicos alargados extraídos con un umbral fijo aparecen fragmentados (Sección~\ref{sec:datos}). El entrenamiento del sintetizador sí usa máscaras umbralizadas de metal real, y su diferencia con el cilindro paramétrico es uno de los desplazamientos de dominio entre entrenamiento y síntesis (Sección~\ref{sec:sintetizador}). Ningún objetivo aísla este aporte, porque ningún brazo de comparación coloca geometrías extraídas por umbral (Sección~\ref{sec:geometria}).

El segundo elemento es un muestreador cuya distribución de poses se compara con la referencia clínica, y el tercero, la codificación multiventana usada para generar, junto con la banda $B_{\delta}$. En las fuentes revisadas no hay un precedente que fije con un valor numérico una banda de generación fuera de la máscara del implante. El uso de varias ventanas de HU proviene de trabajos de MAR, que las aplican a etapas de reconstrucción o a la función de pérdida y no a la codificación de la entrada \cite{wang2025adaptiveweighting,li2024}. El aporte de este trabajo se limita a emplear esas ventanas como codificación de la entrada para generar. Del sintetizador se reclama la combinación del objeto metálico rígido, la codificación multiventana, la banda $B_{\delta}$ y la copia exacta del volumen fuera de $G$, y no cada pieza por separado (Sección~\ref{sec:ea-brecha}). Ningún objetivo aísla la codificación ni la banda, porque el Objetivo~3 evalúa el sintetizador completo.

La justificación metodológica tiene dos partes: el muestreador no se construye con los datos contra los que se compara, y el veredicto negativo del Objetivo~1 aporta un hallazgo propio. La distribución de poses del muestreador se fijó por escrito en su preinscripción, antes de calcular cualquier distancia contra la referencia clínica, y ningún parámetro del muestreador se toma de ella. El hallazgo del Objetivo~1 es que ninguna de las seis combinaciones bajó de 25~HU, y que en el autoencoder entrenado con TC la sola saturación a su rango de salida ya superaba ese criterio (Sección~\ref{sec:obj1}).

\introsection{Alcance y limitaciones / restricciones}

El alcance queda acotado en cuatro puntos: la cohorte, el implante, el nivel sacro y la forma de la síntesis. La anatomía receptora y las observaciones reales de artefacto provienen del subconjunto local de CTPelvic1K, 178 de los 1\,184 volúmenes que declara la publicación \cite{liu2021ctpelvic1k}. El único implante que se coloca y se sintetiza es un tornillo transilíaco-transsacro paramétrico y rígido; los implantes reales del conjunto no aportan geometría. Es el único implante con un corredor óseo medible sobre cada volumen y con una distribución clínica publicada de grados de brecha, y en las fuentes revisadas las placas y los demás tornillos no tienen ninguna de las dos. Esa distribución, como el marco de Kaiser et al., se refiere a tornillos iliosacros, y entra como referencia externa de grados y no como prueba de que ambos tornillos sean el mismo (Sección~\ref{sec:mt-iliosacra}).

En el nivel sacro, la comparación con la referencia clínica se limita a S1. El segundo corredor bajo S1, un corredor óseo cuyo nivel vertebral varía entre volúmenes, queda fuera de esa comparación, porque ninguna de las fuentes revisadas reporta para él una distribución ordinal clínica. Su eje está medido, pero el muestreo de poses en él no se ha ejecutado \GAPDATO{preinscripción y corrida del muestreo de poses en el segundo corredor bajo S1}. La síntesis opera sobre parches alrededor del implante, y dentro de cada parche solo genera la región $G$; fuera de ella copia el volumen de origen (Sección~\ref{sec:sintetizador}).

Quedan fuera de alcance cuatro tareas, cada una por una razón distinta:
\begin{itemize}
    \item \textbf{Evaluación de segmentación posterior.} Medir el efecto de los implantes sintéticos sobre una red de segmentación ósea con el coeficiente de Dice y la distancia de Hausdorff al percentil~95 (HD95) exigiría casos reales con metal y anotación ósea. En CTPelvic1K son 14 (véase la justificación), y no está verificado que sus anotaciones estén disponibles localmente. Además, localmente solo se dispone de una parte de la colección. Un resultado de segmentación descansaría en una cohorte pequeña y parcialmente anotada.
    \item \textbf{Ablaciones del muestreador.} Retirar uno a uno los componentes del muestreador vigente, el anclaje al eje del corredor y la escala de la perturbación, para medir el efecto de cada uno se difiere a trabajo futuro, para concentrar el tiempo en la compuerta y el muestreador.
    \item \textbf{Reimplementación validada de XCIST/CatSim.} La validación publicada del simulador es preliminar y no incluye un estudio de artefacto metálico \cite{wu2022xcist}, de modo que una reimplementación propia no tendría ninguna validación en metal. En su lugar se adopta el protocolo físico de Peters et al.~\cite{peters2025hybrid}, que se ejecuta sobre ese mismo simulador, publica sus métricas y valida su simulación de artefactos metálicos contra un fantoma físico. Esa validación tampoco cubre la síntesis (véanse las limitaciones).
    \item \textbf{Estratificación por dismorfismo sacro.} Gardner et al.~\cite{gardner2010safezones} clasifican el sacro como dismórfico sobre radiografía simple, por rasgos como el ala sacra angulada y ascendente o los forámenes neurales superiores no circulares, sin exigir que estén todos presentes. El efecto del dismorfismo sobre el tamaño de la zona segura no es consistente entre estudios. Esos autores reportan en S1 un área transversal mínima de la zona segura de 222~mm$^2$ en sacros dismórficos frente a 346~mm$^2$ en normales, pero recogen un estudio previo que no halló diferencia. Por eso el muestreador no estratifica por dismorfismo y mide, en cambio, la geometría del corredor sobre cada volumen.
\end{itemize}

El diseño descansa en supuestos, que son afirmaciones no verificadas sobre la física del artefacto o sobre las fuentes, y en convenciones, que son valores fijados por este trabajo; la Sección~\ref{sec:amenazas} discute su efecto sobre los resultados. El primer supuesto es que la apariencia del artefacto metálico puede generarse en el dominio de imagen, sin datos de proyección, aunque De Man et al.~\cite{deman1999} caracterizan sus mecanismos como no lineales y los aíslan modificando el sinograma. La comparación con el protocolo físico es la que puede contradecirlo (véanse las limitaciones). El segundo supuesto es que la serie navegada de la referencia clínica, cuyo nivel sacro Zwingmann et al. no declaran, corresponde a S1. El tercero es que la relación entre la máscara y el artefacto no depende del tipo de implante, lo que permite entrenar el sintetizador con componentes metálicos de todo tipo.

Además de los valores declarados en la preinscripción del muestreador (Tabla~\ref{tab:preinscripcion}), el diseño fija otras convenciones. Ninguna fuente revisada calibra el ancho de unos 12~mm de $B_{\delta}$. Los límites de grado de la escala de brecha cortical, de 2~mm de ancho, provienen de la literatura de tornillos pediculares \cite{gertzbein1990,mirza2003} y llegan a la fijación iliosacra a través de Smith et al.~\cite{smith2006iliosacral}. Su traslado al corredor iliosacro es una convención geométrica y un supuesto explícito, no un umbral clínico validado. Los cuatro grados de brecha se tratan como equiespaciados al calcular la distancia de Wasserstein-1. La desviación estándar de las perturbaciones del eje se fija en la mitad del margen cortical de longitud útil de 5~mm que exigen Kaiser et al.~\cite{kaiser2014dysmorphism}. Esa convención no tiene fuente y se justifica por su consecuencia: con perturbaciones de magnitud seminormal, cerca del 95\,\% de las magnitudes quedan dentro de ese margen (Sección~\ref{sec:poses}).

Las limitaciones se resumen aquí por componente, y la Sección~\ref{sec:amenazas} las discute en detalle. En la colocación, la referencia clínica proviene de pelvis fracturadas, mientras que las pelvis receptoras se eligieron sin osteosíntesis y no se verificaron libres de fractura. Una malreducción residual en la referencia clínica estrecharía los corredores de sus pacientes de una forma que las pelvis receptoras no reproducen, y una fractura no detectada en estas los estrecharía por una causa distinta de la variabilidad anatómica. En un corredor más estrecho que el tornillo con que se mide la brecha, el propio eje ya perfora y el grado~0 es inalcanzable por anatomía (Sección~\ref{sec:sap}). Por eso ambas diferencias pueden mover la distancia de Wasserstein-1 sin que intervenga el muestreador. Un cribado ciego de fractura sobre 30 casos de la cohorte del Objetivo~2, 15 de corredor estrecho y 15 de control, halló fractura confirmada en 20 de los 30, 10 en cada grupo (Sección~\ref{sec:datos}). Lo hizo un médico sin especialidad, de modo que es un cribado y no una lectura de especialista, y muestra que la cohorte receptora no es una colección de pelvis sanas.

Las cifras principales del Objetivo~2 dependen de una variante de la segmentación anatómica cuya elección quedó condicionada a revisar 16 casos, marcados por diferencia de diámetro entre variantes o por desplazamiento de la caja de una estructura. La autora, con apoyo de un médico sin especialidad, revisó los 16 sobre láminas: 14 salieron conformes y 2 fallaron con las dos variantes (Sección~\ref{sec:corredor}). Ninguno falló solo con la variante principal, que es el único veredicto que habría reabierto su elección, y la variante se mantiene. La publicación de TotalSegmentator, la herramienta de esa segmentación, no reporta exactitud para la etiqueta de S1 \cite{wasserthal2023}, de modo que esa exactitud no se asume y un control verifica el nivel S1 de cada caso. Entre los candidatos a la cohorte del Objetivo~2, la discordancia de nivel excluyó a 11 de 34 pacientes con objetos no ortopédicos y a 7 de 57 sin objeto. Si el corredor difiriera entre esos grupos, la distribución de grados heredaría la composición de la cohorte. La brecha se mide contra una envolvente ósea, una máscara de hueso cuyo cierre morfológico podría ocupar el canal sacro y los forámenes, y en ese caso una perforación hacia ellos daría brecha nula \GAPDATO{comprobación de que la envolvente ósea deja fuera el canal sacro y los forámenes sacros}.

En la apariencia, las observaciones reales de artefacto tienen una reconstrucción no reportada, que pudo incluir una MAR del fabricante o una reconstrucción monoenergética, y ambas cambian la apariencia del artefacto \cite{selles2024marreview}. Por eso sirven para comparar la apariencia y no como verdad física. El protocolo físico adoptado se validó en dos dimensiones y para MAR, sin cubrir el paso híbrido con imágenes clínicas, de modo que usarlo como comparación para la síntesis exige una adaptación que no hereda esa validación. Como $B_{\delta}$ trunca por diseño las rayas lejanas y el protocolo físico no, la comparación entre ambos no puede mostrar si el sintetizador reproduciría las rayas más allá de la banda.

Entre el entrenamiento del sintetizador y la síntesis hay tres desplazamientos de dominio, que limitan la transferencia de lo aprendido sobre metal real a los tornillos paramétricos. Las máscaras de entrenamiento se umbralizan sobre metal real, y la de síntesis es el cilindro paramétrico liso. El contexto de entrenamiento trae rayas lejanas que no tienen los volúmenes sin metal que reciben la síntesis. Los parches que solo contienen banda son menos frecuentes en el entrenamiento que en la síntesis. El primero se cuantifica con la fragmentación y el diámetro de los tornillos medidos en la auditoría, el segundo se mide en la evaluación del sintetizador y el tercero, con la razón entre tipos de parche (Sección~\ref{sec:sintetizador}). Por último, el Objetivo~3 no tiene todavía resultados \GAPDATO{el sintetizador generó muestras sobre componentes metálicos reales y la cadena completa se ejecutó sobre pacientes de validación, pero no existe aún resultado del Objetivo~3 porque el modelo final no está elegido}, y su comparación de equivalencia frente al protocolo físico depende del plazo \GAPDEC{si la comparación de equivalencia frente al protocolo físico se mantiene o pasa a trabajo futuro, contingencia de plazo registrada por la autora}.
   % Bloque (a)
\chapter{MARCO TEÓRICO}\label{cap:marco}
% Bloque (b), parte teorica. Ver redaccion/MAPA.md.
% Pautas de la plantilla UTEC movidas a redaccion/RUBRICA.md (P-MT).
% Fuentes: docs/03-glosario.md; fichas de docs/literatura/ (deman1999, glover1980nonlinear,
% park2015ct, selles2024marreview, deman2007catsim, wu2022xcist, ho2020denoising, song2021ddim,
% nichol2021improved, rombach2022latentdiffusion, lugmayr2022repaint, zhang2025diffboost,
% dorjsembe2024, zhang2023controlnet, kaiser2014dysmorphism, mclaren2021corridor,
% gardner2010safezones, smith2006iliosacral, zwingmann2009navigated, routt1997, wasserthal2023);
% tesis/main.tex. Fuente de cada afirmacion: redaccion/rondas/capitulo1-r00-respuesta.md.

Este capítulo reúne los conceptos que el Capítulo~\ref{cap:propuesta} usa sin volver a definirlos. Está escrito para un lector de computación que no conoce la física de la tomografía ni la cirugía de la pelvis. De cada concepto se da la definición, la notación y la fuente, y se dice en qué pieza de la propuesta interviene. El capítulo no compara trabajos entre sí, que es la función del Capítulo~\ref{cap:estado-arte}, ni describe el diseño, que es la del Capítulo~\ref{cap:propuesta}.

Las secciones van de la física de la imagen a la cirugía, al modelo generativo y a la estadística: el artefacto se explica con los conceptos de la tomografía, y el modelo 2.5D, con los del artefacto. La cirugía va antes del modelo generativo porque el muestreador, que usa sus conceptos, precede en la cadena al sintetizador (Figura~\ref{fig:pipeline}). La Sección~\ref{sec:mt-tc} define qué mide una tomografía y en qué unidades, de lo que dependen la compuerta del Objetivo~1 y la representación del sintetizador. La Sección~\ref{sec:mt-artefactos} explica cómo produce el metal su artefacto y por qué ese artefacto no queda dentro del implante. La Sección~\ref{sec:mt-iliosacra} describe la fijación iliosacra, de la que salen el corredor óseo, la escala de brecha cortical y la referencia clínica del Objetivo~2. La Sección~\ref{sec:mt-difusion} presenta los modelos de difusión con que se formula el sintetizador del Objetivo~3, y la Sección~\ref{sec:mt-estadistica}, el análisis estadístico con que se leen los resultados. La Figura~\ref{fig:mt-conceptos} relaciona cada grupo de conceptos con la pieza de la propuesta que lo usa.

\begin{figure}[htbp]
\centering
\small
\begin{tabular}{p{0.44\linewidth} c p{0.40\linewidth}}
\fbox{\parbox{0.95\linewidth}{\centering Atenuación, sinograma, reconstrucción, unidades Hounsfield, ventanas y codificación multiventana (Sección~\ref{sec:mt-tc})}} & $\rightarrow$ & \fbox{\parbox{0.93\linewidth}{\centering Compuerta de representación (Objetivo~1); canales de entrada y salida del sintetizador (Objetivo~3)}} \\ \noalign{\vspace{8pt}}
\fbox{\parbox{0.95\linewidth}{\centering Mecanismos del artefacto metálico, dominio en que se definen y métricas con que se cuantifica (Sección~\ref{sec:mt-artefactos})}} & $\rightarrow$ & \fbox{\parbox{0.93\linewidth}{\centering Banda de generación extendida $B_{\delta}$ y supuesto del dominio de imagen (Objetivo~3); métricas de apariencia (Objetivos~3 y~4); protocolo físico de comparación}} \\ \noalign{\vspace{8pt}}
\fbox{\parbox{0.95\linewidth}{\centering Tornillo iliosacro, zona segura, corredor óseo, marco de referencia y escala de brecha cortical (Sección~\ref{sec:mt-iliosacra})}} & $\rightarrow$ & \fbox{\parbox{0.93\linewidth}{\centering Muestreador de colocación (Objetivo~2) y SAP (Objetivo~4)}} \\ \noalign{\vspace{8pt}}
\fbox{\parbox{0.95\linewidth}{\centering Procesos directo e inverso, condicionamiento, difusión latente, \emph{inpainting} y modelo 2.5D (Sección~\ref{sec:mt-difusion})}} & $\rightarrow$ & \fbox{\parbox{0.93\linewidth}{\centering Autoencoder examinado en la compuerta (Objetivo~1); sintetizador (Objetivo~3)}} \\ \noalign{\vspace{8pt}}
\fbox{\parbox{0.95\linewidth}{\centering Distancia de Wasserstein-1, superioridad y equivalencia, intervalos por remuestreo (Sección~\ref{sec:mt-estadistica})}} & $\rightarrow$ & \fbox{\parbox{0.93\linewidth}{\centering Lectura de los resultados de los Objetivos~1 a~3; definición de la distancia en SAP (Objetivo~4)}} \\
\end{tabular}
\caption{Conceptos del capítulo y pieza de la propuesta que los usa.}
\label{fig:mt-conceptos}
\end{figure}

\section{Tomografía computarizada y unidades Hounsfield}\label{sec:mt-tc}

Una TC no fotografía el cuerpo: mide cuánto se atenúan los rayos X a lo largo de muchas trayectorias y reconstruye, a partir de esas medidas, un mapa de atenuación. Cada lectura del detector corresponde a un rayo de una vista, y el conjunto de lecturas ordenadas por índice forma el \textbf{sinograma} \cite{deman2007catsim}. En el modelo de adquisición de CatSim, la señal de cada índice del sinograma suma, sobre las energías del espectro, los fotones que llegarían al detector sin objeto, atenuados según la longitud que el rayo recorre en cada material. Esa atenuación depende del coeficiente de atenuación lineal de cada material a cada energía \cite{deman2007catsim}. Abadi et al.~\cite{abadi2019}, en el simulador DukeSim, y Wang et al.~\cite{wang2019cochlear}, en su simulación de implantes cocleares, calculan esa atenuación con la ley de Beer-Lambert; Wang et al. la discretizan en cinco energías.

La imagen se obtiene invirtiendo ese proceso de medida. Las simulaciones de De Man et al.~\cite{deman1999,deman2007catsim} y de Karageorgos et al.~\cite{karageorgos2024ddpm} reconstruyen con retroproyección filtrada (\emph{filtered back-projection}). Park et al.~\cite{park2015ct} señalan que ese algoritmo supone un modelo monocromático de las proyecciones, el de Radon, y que las proyecciones de un haz policromático se apartan de él de forma no lineal. Esa discrepancia reaparece en la imagen como artefacto (Sección~\ref{sec:mt-artefactos}). Este documento distingue dos dominios: el \textbf{dominio de proyección}, el del sinograma, y el \textbf{dominio de imagen}, el de los valores reconstruidos por vóxel. La cohorte de este trabajo solo contiene volúmenes ya reconstruidos (Sección~\ref{sec:datos}), de modo que lo que este trabajo mide, y lo que genera el sintetizador, ocurre en el dominio de imagen. Solo el protocolo físico de comparación, previsto en la Sección~\ref{sec:apariencia}, pasa por proyecciones simuladas.

Los valores de TC se expresan en \textbf{unidades Hounsfield} (HU), una escala normalizada de la atenuación reconstruida en la que el agua vale 0~HU y el aire $-1000$~HU por definición \cite{wu2022xcist}. La conversión entre coeficiente de atenuación y HU se apoya en el coeficiente de absorción del agua en función de la energía \cite{wang2019cochlear} \GAPLIT{fuente de referencia de física de TC que enuncie la ley de Beer-Lambert con el coeficiente de atenuación lineal, dé la definición de la escala Hounsfield como transformación de ese coeficiente y los fundamentos de la reconstrucción por retroproyección filtrada}. El protocolo de evaluación adoptado define el hueso como los vóxeles sobre 150~HU fuera de la geometría metálica \cite{peters2025hybrid}. Dos fuentes de MAR segmentan el metal clínico con 2500~HU \cite{wang2025adaptiveweighting,li2024}, umbral con que este trabajo criba la cohorte y delimita el metal clínico de sus volúmenes (Secciones~\ref{sec:datos} y~\ref{sec:sintetizador}). Los dos umbrales son definiciones operativas de las fuentes que los usan y no límites físicos del tejido. En la cohorte primaria, con el recorte por defecto, la mediana de la fracción de vóxeles a 150~HU o menos dentro de las máscaras sacras fue 0.42 a lo largo del eje del corredor más ancho. Por eso el corredor se mide sobre máscaras anatómicas y no con el umbral de 150~HU, que dejaría fuera esa fracción (Sección~\ref{sec:corredor}).

Varias de las redes de síntesis y de MAR revisadas normalizan los HU con una ventana antes de procesarlos (Sección~\ref{sec:ea-representacion}). Una \textbf{ventana} es un intervalo $[w_{\min}, w_{\max}]$ de HU que se satura y se lleva linealmente a $[0, 1]$ \cite{wang2025adaptiveweighting}:
\begin{equation}
\tilde{x} = \frac{\min\bigl(\max(x, w_{\min}), w_{\max}\bigr) - w_{\min}}{w_{\max} - w_{\min}},
\label{eq:ventana}
\end{equation}
donde $x$ es el valor de TC de un vóxel y $\tilde{x}$ su valor normalizado. Lo que cae fuera de la ventana se \textbf{satura}: todos los valores por encima de $w_{\max}$ se vuelven indistinguibles entre sí, y lo mismo ocurre bajo $w_{\min}$. Dentro de la ventana, un mismo intervalo de HU ocupa una fracción del rango normalizado inversamente proporcional al ancho de la ventana. Por eso una ventana ancha conserva más rango pero comprime los contrastes pequeños del tejido, y una estrecha los separa pero satura más valores, entre ellos los del hueso denso y el metal.

La \textbf{codificación multiventana} usa varias ventanas a la vez, una por canal, y evita así elegir entre conservar el rango y separar los contrastes. En la MAR, Wang et al.~\cite{wang2025adaptiveweighting} trabajan con tres ventanas, $[-1000, 2000]$, $[-320, 480]$ y $[-160, 240]$~HU, aunque en una cascada de etapas y no como canales de entrada (Sección~\ref{sec:ea-representacion}). El paso de ida y vuelta lleva los HU a la representación y de vuelta a HU, y su error es la diferencia, en HU, entre el valor original de cada vóxel y el recuperado.

Este trabajo usa la codificación multiventana en dos piezas: la compuerta del Objetivo~1 y el sintetizador del Objetivo~3. El Objetivo~1 examinó tres codificaciones (Sección~\ref{sec:obj1}). La primera usa como canales las ventanas de Wang et al., cuyo techo de 2000~HU queda bajo el umbral de 2500~HU del metal, de modo que el metal satura también en la ventana ancha. La segunda eleva ese techo a 20\,000~HU, y la tercera sustituye la ventana ancha por una compresión arcoseno hiperbólico sobre $[-1000, 20\,000]$~HU, que a diferencia de la Ecuación~\ref{eq:ventana} no es lineal \GAPDATO{forma y parámetros de la compresión arcoseno hiperbólico, que solo constan en el código del experimento y no en un registro de sus resultados}. La compuerta mide el error de ida y vuelta como error absoluto medio dentro del hueso (Ecuación~\ref{eq:mae}). El sintetizador usa la codificación multiventana como representación de sus canales de entrada y salida, sin la etapa del autoencoder (Sección~\ref{sec:sintetizador}), y la variante que adopta no se ha fijado \GAPDEC{qué codificación multiventana y qué precisión numérica adopta el sintetizador, pendiente de su preinscripción}.

\section{Artefactos inducidos por metal}\label{sec:mt-artefactos}

Un \textbf{artefacto metálico} no es la imagen del metal, sino el error que el metal introduce en la reconstrucción del tejido que lo rodea. Selles et al.~\cite{selles2024marreview} lo describen como rayas claras y oscuras que ocultan estructuras anatómicas. Según la misma revisión, su severidad depende del tamaño, la forma y la aleación del implante, y los materiales de osteosíntesis caen en una categoría de artefacto intermedia, asignada de forma cualitativa y sin criterio numérico.

Las fuentes revisadas atribuyen el artefacto a varios mecanismos físicos; sus listas comparten el endurecimiento del haz, la dispersión y los efectos de borde, y difieren en el resto. Selles et al. nombran como principales el endurecimiento del haz, la inanición de fotones, la dispersión y los efectos de borde. De Man et al.~\cite{deman1999} aíslan cada causa en una simulación bidimensional y concluyen que las más importantes son el endurecimiento del haz, la dispersión, el ruido y el efecto exponencial de gradiente de borde. También concluyen que todas las causas que estudian producen rayas, sin asignar a ninguna un peso numérico.

En la simulación bidimensional de De Man et al., el endurecimiento del haz y la dispersión producen rayas muy parecidas entre sí \cite{deman1999}. El \textbf{endurecimiento del haz} (\emph{beam hardening}) ocurre porque el material absorbe con preferencia los fotones de baja energía, de modo que el espectro que sale del objeto se desplaza hacia energías más altas \cite{selles2024marreview}. Park et al.~\cite{park2015ct} separan dos efectos del endurecimiento: un error de los valores dentro del propio objeto (\emph{cupping}) y rayas que corrompen la imagen fuera de la región que las produce. En esa simulación, el endurecimiento produce rayas oscuras en las direcciones de mayor atenuación y, con dos o más objetos metálicos, rayas oscuras que los conectan \cite{deman1999}. La \textbf{dispersión} (\emph{scatter}) aumenta con la alta densidad electrónica del metal \cite{selles2024marreview}. Para producir rayas apreciables, en la misma simulación basta una razón entre radiación dispersa y primaria de 0.0001, elegida de forma arbitraria \cite{deman1999}.

La \textbf{inanición de fotones} (\emph{photon starvation}) aparece cuando un rayo atraviesa tanto metal que llegan muy pocos fotones al detector, y faltan datos de proyección esenciales \cite{selles2024marreview}. De Man et al. no nombran la inanición de fotones; describen como artefacto de ruido líneas finas, alternadamente oscuras y claras, cuya intensidad depende de la atenuación total integrada a lo largo del rayo \cite{deman1999}. Según ellos, ese artefacto de ruido es no lineal, como los del endurecimiento del haz y la dispersión, y su severidad depende de cuánta dispersión y endurecimiento haya, de modo que en su simulación las causas no actúan por separado.

El \textbf{efecto exponencial de gradiente de borde} (\emph{exponential edge-gradient effect}) produce, según recuerdan De Man et al., rayas tangentes a los bordes rectos largos, y en su simulación se observan rayas que irradian desde el metal \cite{deman1999}. Selles et al. describen los efectos de borde como rayas oscuras o claras alineadas con el borde del metal. Ninguna de las dos descripciones incluye una distancia, y el artículo de De Man et al. no menciona una medida de hasta dónde llegan las rayas de ninguna de las causas que estudia.

A diferencia de los anteriores, el \textbf{volumen parcial no lineal} depende de la tercera dimensión. Glover y Pelc~\cite{glover1980nonlinear} derivan que, cuando la atenuación varía en la dirección axial dentro del espesor del haz, integrar el flujo y tomar su logaritmo vuelve no lineal la medida. En su análisis, una sola discontinuidad produce sobre todo un error local, mientras que dos o más pueden producir rayas de largo alcance que conectan estructuras. Lo estudian en hueso, sin metal y con una fuente monocromática idealizada, y De Man et al. lo dejan fuera al simular en dos dimensiones \cite{deman1999}: ninguna de las dos fuentes reúne metal y volumen parcial axial.

De estos mecanismos se siguen dos consecuencias para la síntesis. La primera es que el artefacto no queda dentro del implante: las rayas se extienden por el tejido que rodea al metal \cite{deman1999,park2015ct}. Por eso el sintetizador genera también en la banda de generación extendida $B_{\delta}$ alrededor de la máscara del implante, cuyo ancho es una convención de este trabajo porque las fuentes revisadas no dan hasta dónde llegan las rayas (Sección~\ref{sec:ea-representacion}). La segunda es que los mecanismos se definen sobre las proyecciones: De Man et al. aíslan cada causa modificando el sinograma y reconstruyendo después. Como la cohorte solo contiene volúmenes reconstruidos (Sección~\ref{sec:mt-tc}), este trabajo asume que la apariencia del artefacto puede generarse en el dominio de imagen; la Sección~\ref{sec:amenazas} discute ese supuesto.

La MAR reduce el artefacto de una TC ya adquirida, y la simulación física lo produce reproduciendo la adquisición; el sintetizador de este trabajo no hace ni lo uno ni lo otro. CatSim, el simulador del conjunto XCIST \cite{wu2022xcist} sobre el que se ejecuta el protocolo físico de Peters et al., modela el espectro policromático, el ruido cuántico y el electrónico, el volumen parcial no lineal y la dispersión \cite{deman2007catsim}. Este trabajo genera la apariencia del artefacto con un modelo aprendido, sin simular proyecciones, y prevé usar la simulación física como comparación (Sección~\ref{sec:apariencia}).

Para cuantificar el artefacto, este trabajo adopta las métricas del protocolo de Peters et al.~\cite{peters2025hybrid}, con sus nombres publicados. La \emph{streak amplitude} mide las rayas en regiones perpendiculares a ellas, como la diferencia entre el promedio del 5\,\% de desviaciones más altas respecto de la imagen sin metal del mismo caso y el del 5\,\% más bajas. La \emph{bone integrity} compara el hueso, delimitado con el umbral de 150~HU de la Sección~\ref{sec:mt-tc}, mediante el cambio de volumen y el coeficiente de Sørensen-Dice. La \emph{metal integrity} compara de forma análoga la máscara del metal, con un umbral que se adapta a cada región. Las tres se diseñaron para evaluar MAR, y su uso para evaluar síntesis exige invertirlas; la definición operativa de esa inversión está pendiente para cada métrica (Sección~\ref{sec:apariencia}). El Objetivo~3 toma la \emph{streak amplitude} como su único criterio primario.

\section{Fijación iliosacra y corredores óseos}\label{sec:mt-iliosacra}

La fijación iliosacra trata fracturas inestables del anillo pélvico posterior. Zwingmann et al.~\cite{zwingmann2009navigated} incluyen en su estudio solo fracturas de tipos~B y~C de la clasificación de Tile y Pennal. La de tipo~B es inestable para la rotación, y la de tipo~C, para la rotación y el desplazamiento vertical. En ese estudio, los tornillos se colocan por vía percutánea, y la distribución de sus grados de brecha forma la referencia clínica del Objetivo~2.

En Smith et al.~\cite{smith2006iliosacral}, el \textbf{tornillo iliosacro} entra por el ilion y llega a la primera o a la segunda vértebra sacra, de modo que atraviese tres corticales, dos del ilion y una del ala sacra. El ala sacra desciende a cada lado, hacia lateral y caudal, desde el cuerpo vertebral superior del sacro \cite{routt1997} \GAPLIT{fuente de anatomía de la pelvis que defina e ilustre la cortical ósea, el platillo vertebral, el ala sacra, el foramen neural, la tabla externa del ilion y la articulación sacroilíaca}. Un tornillo transiliosacro, en cambio, cruza las dos articulaciones sacroilíacas y sale por la tabla externa del ilion opuesto \cite{mclaren2021corridor}. Este trabajo trata como no intercambiables las medidas de corredor de los dos tipos: las de McLaren et al. son de tornillos transiliosacros, y Kaiser et al.~\cite{kaiser2014dysmorphism} eligen su umbral de 10~mm para el paso de un tornillo iliosacro. La Sección~\ref{sec:sap}, sin embargo, trata el tornillo de este trabajo como uno que entra y sale por la cortical del ilion \GAPDEC{qué tornillo representa el corredor que mide este trabajo, que el registro describe de la cortical externa de un ilion a la del ilion opuesto, y con ello si le aplican el umbral de 10~mm y la holgura radial de Kaiser et al., definidos para el tornillo iliosacro, y si es del mismo tipo que los tornillos de la referencia clínica, que Zwingmann et al. llaman iliosacros y transiliosacros}.

La colocación percutánea se guía con imagen intraoperatoria. Routt et al.~\cite{routt1997} describen la técnica con el paciente en decúbito supino y fluoroscopia en tres planos. Zwingmann et al. comparan la guía fluoroscópica, que llaman convencional, con navegación computarizada sobre una adquisición tridimensional intraoperatoria hecha con un equipo que rota $190^{\circ}$ alrededor del campo quirúrgico. De esa comparación salen la serie navegada y la serie convencional. En ambas, la posición final de cada tornillo se evaluó sobre una TC posoperatoria.

El tornillo debe quedar dentro del hueso en una región del sacro superior, la \textbf{zona segura}. Routt et al. la describen en el ala sacra, limitada por delante y por arriba por la cortical del ala, y por detrás y por abajo por el túnel del foramen neural. Cerca de ella quedan la raíz nerviosa de L5 sobre el ala, el canal espinal y los vasos ilíacos por delante del ala, y un tornillo que sale de la zona se acerca a esas estructuras \cite{routt1997}. Gardner et al.~\cite{gardner2010safezones} representan la zona como dos conos unidos por la punta y miden su sección transversal mínima donde los conos se encuentran, en un plano ortogonal al eje de la zona. Este trabajo no usa la zona segura como variable: la presenta como antecedente del corredor óseo, que es lo que el muestreador mide (Sección~\ref{sec:muestreador}).

El \textbf{corredor óseo} extiende esa idea a toda la trayectoria: es la región de hueso por la que puede pasar un tornillo sin atravesar la cortical. McLaren et al.~\cite{mclaren2021corridor} lo resumen en un solo número por segmento sacro. Colocan una recta dentro de los contornos corticales y aumentan su diámetro hasta que toca y atraviesa la cortical en al menos tres puntos, lo que da el diámetro máximo de la trayectoria. Gardner et al., Kaiser et al.~\cite{kaiser2014dysmorphism} y McLaren et al. emplean un umbral de 10~mm sobre el tamaño del corredor. Los tres lo toman de trabajos anteriores, y McLaren et al. declaran además que el corredor mínimo no está establecido. Kaiser et al. lo justifican como una holgura de 1 a 2~mm alrededor de un tornillo de 6.3 a 8~mm, que la Sección~\ref{sec:corredor} lee como holgura radial por lado. Esa sección describe también cómo este trabajo mide el corredor, sobre máscaras anatómicas de TotalSegmentator \cite{wasserthal2023}, y qué criterio de viabilidad adopta en lugar de ese umbral.

Expresar la pose de un tornillo exige un marco de referencia anatómico que se pueda calcular sobre la TC. Kaiser et al. reformatean el volumen a lo largo del eje del sacro, perpendicular al platillo superior de S1. Miden la angulación del corredor contra la línea que une las crestas ilíacas y contra la que une las espinas ilíacas posteriores. El marco incluye además una regla de longitud útil, que exige al menos 5~mm hasta la cortical a cada lado del eje y que la Sección~\ref{sec:corredor} lee como margen cortical. Este trabajo expresa sus poses en ese marco, y la Sección~\ref{sec:marco-metal} examina si el platillo de S1, las crestas y las espinas ilíacas se pueden localizar en volúmenes con metal.

Smith et al.~\cite{smith2006iliosacral} juzgan la posición de un tornillo frente a una posición ideal. La definen enteramente dentro de los márgenes corticales del sacro, paralela al platillo correspondiente y a media altura entre el platillo superior de S1 y el foramen neural de S1. Distinguen tres tipos de perforación según la cortical que el tornillo atraviesa: la anterior del sacro, la del canal espinal y la de los platillos vertebrales. Smith et al. gradúan esa perforación, que en este documento se llama brecha cortical, en cuatro niveles. El grado~0 es la ausencia de perforación, el grado~1 una perforación menor que 2~mm, el grado~2 una de 2 a 4~mm y el grado~3 una mayor que 4~mm. En este trabajo, el grado se calcula a partir de la protrusión del implante fuera de una envolvente ósea segmentada, que aproxima esa perforación (Secciones~\ref{sec:sap} y~\ref{sec:amenazas}).

Las fuentes revisadas no juzgan la colocación con una misma escala. Smith et al. declaran que toman su escala de la clasificación de tornillos pediculares, y le suman una escala angular \GAPDEC{si SAP incorpora la dimensión angular de la escala de Smith et al. o declara que usa solo la de perforación, y con qué justificación}. Zwingmann et al. aplican la escala de perforación (Sección~\ref{sec:sap}). Hinsche et al.~\cite{hinsche2002fluoroscopy} usan, en cambio, una definición binaria: consideran insegura una colocación cuya trayectoria perfora una de las corticales comprometiendo estructuras neurovasculares, sin grados de profundidad.

La escala es ordinal: sus grados están ordenados por la profundidad de la perforación, pero no cubren intervalos del mismo ancho. Los grados~1 y~2 cubren 2~mm cada uno, y el grado~3 no tiene límite superior. Tratar los grados como equiespaciados, como hace la distancia de la Sección~\ref{sec:mt-estadistica}, es entonces una convención y no una propiedad de la escala (Sección~\ref{sec:sap}).

La escala de brecha y la viabilidad del corredor son dos de los tres componentes de SAP; el tercero es la fracción por zona de densidad. La motivación de esa fracción es la heterogeneidad de la masa ósea pélvica que describen Arand et al.~\cite{arand2019pelvicring} con un modelo medio de los valores de gris del anillo pélvico. En ese modelo, los valores del ala sacra son más bajos que los del cuerpo de S1 (Sección~\ref{sec:poses}). SAP reporta esa fracción de forma descriptiva, sin compararla con la referencia clínica, y su definición operativa está pendiente \GAPDEC{fuente y definición operativa de la fracción por zona de densidad de SAP}.

\section{Modelos de difusión y síntesis condicionada}\label{sec:mt-difusion}

Un modelo generativo aprende, a partir de ejemplos, una distribución de la que extrae muestras nuevas. Los modelos de difusión son una familia de ellos, y Kazerouni et al.~\cite{kazerouni2023diffusionsurvey} revisan su uso en imagen médica. Ho et al.~\cite{ho2020denoising} formulan la generación como una cadena de Markov que aprende a invertir un \textbf{proceso directo} que añade ruido gaussiano en $n_{\max}$ pasos; en sus experimentos, $n_{\max} = 1000$. Dorjsembe et al.~\cite{dorjsembe2024} resumen el propósito de ese proceso directo: borrar por completo la estructura original de los datos. La imagen ruidosa del paso $n$ se obtiene de la original en una sola operación, como la escriben Zhang et al.~\cite{zhang2025diffboost} en su método DiffBoost y Dorjsembe et al.:
\begin{equation}
\mathbf{x}_n = \sqrt{\bar{\gamma}_n}\,\mathbf{x}_0 + \sqrt{1 - \bar{\gamma}_n}\,\boldsymbol{\xi}, \qquad \boldsymbol{\xi} \sim \mathcal{N}(\mathbf{0}, \mathbf{I}),
\label{eq:difusion-directa}
\end{equation}
donde $\mathbf{x}_0$ es la imagen original, $\boldsymbol{\xi}$ un ruido gaussiano estándar del mismo tamaño y $\bar{\gamma}_n$ un coeficiente que fija el \textbf{calendario de ruido}. Las fuentes escriben ese coeficiente como $\bar{\alpha}_t$; aquí se cambia de símbolo porque $\alpha$ y $t$ designan otras magnitudes en el Capítulo~\ref{cap:propuesta}. Nichol y Dhariwal~\cite{nichol2021improved} definen un calendario coseno para ese coeficiente y obtienen, de sus cocientes sucesivos, las varianzas $\beta_n$ de cada paso; Zhang et al.~\cite{zhang2025diffboost} las sitúan en el intervalo $(0, 1)$.

En la formulación de Ho et al., la red no predice la imagen limpia, sino el ruido añadido, con un objetivo simplificado que Rombach et al.~\cite{rombach2022latentdiffusion} escriben como
\begin{equation}
\mathcal{L} = \mathbb{E}_{\mathbf{x}_0,\, \boldsymbol{\xi},\, n}\Bigl[\bigl\lVert \boldsymbol{\xi} - \boldsymbol{\xi}_{\theta}(\mathbf{x}_n, n) \bigr\rVert_2^2\Bigr],
\label{eq:perdida-difusion}
\end{equation}
con $n$ tomado de forma uniforme entre 1 y $n_{\max}$, donde $\boldsymbol{\xi}_{\theta}$ es la red con parámetros $\theta$. Para generar, el \textbf{proceso inverso} parte de ruido puro y aplica la red paso a paso, del último al primero. En Ho et al. y en Song et al.~\cite{song2021ddim}, esa red es una U-Net. Ronneberger et al.~\cite{ronnenberger2015unet} la definen como una red completamente convolucional con una ruta que contrae la resolución y otra simétrica que la expande, y que combina las características de alta resolución de la primera con la segunda. El borrador del diseño del sintetizador propone también una U-Net, arquitectura que aún no se ha fijado (Sección~\ref{sec:sintetizador}).

El costo de generar crece con el número de pasos del proceso inverso. Song et al. definen una familia de procesos no markovianos que conserva el objetivo de entrenamiento de Ho et al. y permite muestrear de forma determinista y con menos pasos, sin reentrenar la red. Reportan una generación entre 10 y 50 veces más rápida, en tiempo de reloj, que la del muestreo de Ho et al. Con ello, en un modelo entrenado con el objetivo de Ho et al., el \textbf{procedimiento de muestreo} se elige por separado del entrenamiento. Las tres fuentes evalúan imágenes naturales y ninguna trabaja con TC ni con metal \cite{ho2020denoising,song2021ddim,nichol2021improved}, de modo que la aceleración que reportan Song et al. no está medida sobre TC. El borrador del diseño del sintetizador propone el muestreo determinista de Song et al. y no menciona el calendario de ruido; ambos quedan pendientes, con el resto de su entrenamiento y muestreo (Sección~\ref{sec:sintetizador}).

Un modelo condicionado genera a partir de una condición, además del ruido. Rombach et al. distinguen dos vías: concatenar a la entrada de la red una información espacialmente alineada con la imagen, como una máscara, o introducir condiciones no espaciales, como un texto, mediante atención cruzada. Dorjsembe et al. siguen la primera vía y concatenan la máscara como canal adicional en cada paso. Zhang et al.~\cite{zhang2023controlnet} proponen ControlNet, que congela un modelo de difusión preentrenado y entrena una copia de su codificador que recibe la condición espacial. Esa copia se une al modelo congelado mediante convoluciones inicializadas en cero.

La difusión latente separa la generación en dos etapas \cite{rombach2022latentdiffusion}: un autoencoder comprime la imagen a un espacio latente de menor resolución, y un modelo de difusión opera en ese espacio. Al generar, el decodificador del autoencoder devuelve la imagen en una sola pasada. Rombach et al. entrenan el autoencoder con una pérdida perceptual y una adversarial, para evitar el desenfoque que producen las pérdidas por píxel, y describen su compresión como una etapa que elimina el detalle de alta frecuencia. También advierten que la capacidad de reconstrucción del autoencoder puede volverse un cuello de botella cuando la aplicación exige exactitud fina por píxel. Este trabajo lee como una aplicación de ese tipo conservar los HU del hueso con un error absoluto medio menor que 25~HU (Sección~\ref{sec:obj1}). Por eso el Objetivo~1 midió la ida y vuelta por el autoencoder como condición para adoptar la difusión latente, que el veredicto negativo descartó. Como ControlNet presupone un modelo base congelado y su espacio latente, quedó descartado con ella (Sección~\ref{sec:sintetizador}).

El \emph{inpainting} regenera una región enmascarada de la imagen y conserva el resto. En las fuentes revisadas, el \emph{inpainting} con un modelo de difusión sigue una de dos formas. Lugmayr et al.~\cite{lugmayr2022repaint} usan un modelo preentrenado sin condición y alteran solo el muestreo. En cada paso combinan la región conocida, muestreada de la entrada, con la desconocida, generada por el modelo, y además retroceden y avanzan en el tiempo de difusión para armonizar el borde entre ambas. La otra forma es entrenar un modelo condicionado, como el de \emph{inpainting} que Rombach et al. condicionan por concatenación. El sintetizador de este trabajo sigue esta segunda forma: recibe el parche con la región de generación $G$ borrada y las máscaras, genera solo dentro de $G$ y copia el resto del volumen de origen (Sección~\ref{sec:sintetizador}). No se ha fijado si su procedimiento de muestreo toma algo del de Lugmayr et al. o del método LeFusion de Zhang et al.~\cite{zhang2025lefusion} \GAPDEC{descripción y atribución del procedimiento de muestreo de la difusión del sintetizador frente a los de Lugmayr et al. y de LeFusion, pendiente en el registro del proyecto}.

El sintetizador es \textbf{2.5D}: opera sobre cortes axiales, pero recibe como contexto cortes contiguos al que genera, sin ser un modelo tridimensional completo. El borrador de su diseño propone tres cortes contiguos y generar el central \GAPDEC{número de cortes contiguos de la entrada 2.5D del sintetizador, que solo propone el borrador de su diseño, pendiente de preinscripción}. Como el volumen parcial no lineal nace de integrar la atenuación dentro del espesor del haz (Sección~\ref{sec:mt-artefactos}), este trabajo lee que unos pocos cortes vecinos aportan información axial pero no reproducen esa integración. El sintetizador, que opera sobre valores reconstruidos, no contiene entonces ese mecanismo y a lo sumo puede aprender la apariencia que deja en los cortes.

\section{Marco estadístico de la evaluación}\label{sec:mt-estadistica}

Los resultados de este trabajo son de tres tipos, y cada uno se analiza de forma distinta. El Objetivo~1 compara un error medio con un umbral fijado de antemano, el Objetivo~2 compara dos distribuciones de grados ordinales y el Objetivo~3 compara brazos pareados sobre los mismos pacientes. En el Objetivo~1, la unidad de análisis es el paciente: el error se promedia por paciente. En el Objetivo~2, la distribución reúne todas las poses de la cohorte, y la referencia clínica cuenta tornillos y no pacientes, de modo que ninguna de las dos distribuciones tiene al paciente como unidad (Sección~\ref{sec:amenazas}). En el Objetivo~3, las pruebas se parean por paciente, y la forma de agregar antes las poses y las regiones de medición de rayas está pendiente \GAPDEC{cómo se agregan por paciente las poses y las regiones de medición de rayas antes de las pruebas}.

Para comparar distribuciones de grados, este trabajo usa la \textbf{distancia de Wasserstein-1}. Sobre los cuatro grados de brecha, tratados como equiespaciados, es la suma, de un grado a otro, de las diferencias absolutas entre las dos funciones de distribución acumulada (Ecuación~\ref{eq:w1}). Su unidad es el grado. Si una distribución pone toda su masa un grado más arriba que la otra, la distancia vale 1; si una está entera en el grado~0 y la otra en el grado~3, vale 3, el máximo con cuatro grados. A diferencia de comparar solo la fracción de grado~0, la distancia usa la distribución completa, y dos distribuciones con la misma fracción de grado~0 difieren si reparten distinto el resto \GAPLIT{fuente de referencia de la distancia de Wasserstein-1 y de su forma para distribuciones sobre una recta}.

El Objetivo~3 plantea dos preguntas sobre sus brazos pareados. Una prueba de \textbf{superioridad} pregunta si una condición supera a la otra, y una de \textbf{equivalencia}, si la diferencia entre ambas es tan pequeña que se pueden tratar como intercambiables. Una diferencia no significativa no responde la segunda pregunta, porque puede deberse a una muestra pequeña y no a una diferencia pequeña. Por eso la equivalencia solo se concluye cuando el intervalo de confianza del 90\,\% de la diferencia pareada cae entero dentro de un margen $[-\Delta, +\Delta]$ fijado de antemano.

La prueba de superioridad del Objetivo~3 es la de rangos con signo de Wilcoxon, pareada y de una cola. Para la comparabilidad con el protocolo físico, el diseño prevé dos pruebas unilaterales de equivalencia, con un margen definido por la variabilidad entre semillas del sintetizador y entre corridas repetidas del protocolo físico, y aún sin preinscribir (Sección~\ref{sec:apariencia}) \GAPDEC{si la comparación de equivalencia frente al protocolo físico se mantiene o pasa a trabajo futuro, contingencia de plazo registrada por la autora}. Con muestras pequeñas, esa comparación puede no concluir ni equivalencia ni diferencia.

La compuerta del Objetivo~1 acompaña su error con un intervalo de confianza por remuestreo de pacientes. El error que decide es la media, sobre los pacientes, del error de cada paciente. Esa media se vuelve a calcular sobre muestras de pacientes tomadas con reemplazo, y el intervalo del 95\,\% sale de la distribución de esas medias (Sección~\ref{sec:obj1}). Se remuestrean pacientes porque el paciente es la unidad de análisis del Objetivo~1 \GAPLIT{fuentes de referencia de la prueba de rangos con signo de Wilcoxon, de las dos pruebas unilaterales de equivalencia y de su relación con el intervalo de confianza del 90\,\%, y del intervalo de confianza por remuestreo}.

La protección contra ajustar el resultado a los datos depende de lo que se fija antes de verlos y de cuántas oportunidades de aprobar da la regla. La regla de la compuerta se fijó antes de correr la prueba que decide, y la distribución de poses del muestreador y su métrica, antes de calcular cualquier distancia (Secciones~\ref{sec:obj1} y~\ref{sec:poses}). En ambos casos hubo, sin embargo, decisiones posteriores a ver datos: esa regla se fijó conociendo un resultado exploratorio desfavorable que incluía a los pacientes de prueba, y en el Objetivo~2 hubo tres (Sección~\ref{sec:amenazas}). La regla de la compuerta aprueba si lo hace cualquiera de sus seis combinaciones de autoencoder y codificación; esa \textbf{multiplicidad} sin corregir solo puede favorecer un veredicto aprobatorio y, como el veredicto fue negativo, no lo compromete. La compuerta se extendió después al espacio latente del generador de volúmenes de TC de Guo et al.~\cite{guo2025maisi}, con la misma regla y los mismos pacientes de prueba, lo que eleva a siete las combinaciones que pueden aprobar. El veredicto de esa extensión se presenta en el capítulo de resultados y, según la decisión que la autorizó, no cambia el diseño del Objetivo~3. La multiplicidad de siete combinaciones solo podría comprometer ese veredicto si fuera aprobatorio (Sección~\ref{sec:amenazas}).

El error absoluto medio y la raíz del error cuadrático medio, los dos errores por vóxel del documento, no son intercambiables. El error absoluto medio promedia el valor absoluto de las diferencias en HU (Ecuación~\ref{eq:mae}). La raíz del error cuadrático medio promedia sus cuadrados y toma la raíz, de modo que, para los mismos residuos, nunca es menor que el error absoluto medio. De ahí que los errores que reporta la literatura de MAR, expresados como raíz del error cuadrático medio, sirvan al Objetivo~1 como orden de magnitud y no como equivalente de su error absoluto medio (Sección~\ref{sec:ea-representacion}).
      % Bloque (b) teorico: marco teorico
\chapter{ESTADO DEL ARTE Y BRECHA}
% Bloque (b), parte comparativa. Esqueleto propuesto: ver redaccion/MAPA.md.
% Pautas de la plantilla UTEC movidas a redaccion/RUBRICA.md (P-EA).

\section{Síntesis de lesiones mediante difusión}

\section{Simulación física de artefactos metálicos}

\section{Planificación y colocación de tornillos iliosacros}

\section{Representación multiventana en reducción de artefactos}

\section{Comparación crítica y brecha}
      % Bloque (b) comparativo: estado del arte y brecha
\chapter{PROPUESTA Y METODOLOGÍA}\label{cap:propuesta}
% Bloque (c), parte formal y diseno experimental. Ver redaccion/MAPA.md.
% Pautas de la plantilla UTEC movidas a redaccion/RUBRICA.md (P-MM).
% Fuentes: tesis/main.tex (Objectives, Expected Results, Pre-registration, Datasets),
% experiments/objetivo2/preinscripcion_muestreador.md, docs/02-datos.md,
% docs/01-decisiones.md. Fuente de cada cifra: redaccion/rondas/capitulo3-r00-respuesta.md
% capitulo3-r01-respuesta.md a capitulo3-r05-respuesta.md.

Este capítulo describe la contribución y el diseño con que se evalúa. La contribución no es un generador único, sino una cadena de dos componentes que se diseñan y se evalúan por separado. El \textbf{muestreador de colocación} decide dónde y con qué pose tridimensional se ubica un tornillo transilíaco-transsacro (transiliosacro). El \textbf{sintetizador} genera la apariencia del implante y de su artefacto metálico local en la tomografía computarizada (TC). Antes de ambos se ubica una compuerta que decide en qué representación puede operar el sintetizador. El capítulo sigue ese orden: visión general, datos, compuerta de representación, muestreador, sintetizador, protocolo de evaluación y amenazas a la validez.

\section{Visión general}\label{sec:vision}

La propuesta se organiza en cuatro objetivos específicos \GAPDEC{alinear la pregunta de investigación y los objetivos de la introducción con los cuatro objetivos de este capítulo, que hoy no coinciden}. El Objetivo~1 evalúa si una representación de la TC conserva las unidades Hounsfield (HU) del hueso con una tolerancia fijada de antemano. El Objetivo~2 diseña el muestreador de colocación y lo evalúa con la admisibilidad quirúrgica de la colocación (SAP, \emph{Surgical Admissibility of Placement}), que es la única métrica que introduce este trabajo. El Objetivo~3 diseña el sintetizador condicionado. El Objetivo~4 formaliza SAP y adopta, para la apariencia, un protocolo publicado con sus métricas y sus nombres originales, sin renombrarlas como contribución.

La cadena recibe un volumen de TC pélvica sin osteosíntesis y devuelve el mismo volumen con un tornillo sintético y su artefacto. El muestreador mide sobre cada volumen un corredor óseo, perturba su eje y entrega una pose; esa pose, junto con un modelo paramétrico del tornillo, define la \textbf{máscara del implante $M$}. Alrededor de $M$ se añade una \textbf{banda de generación extendida $B_{\delta}$} (\emph{generation band}), y ambas forman la \textbf{región de generación} $G = M \cup B_{\delta}$. El sintetizador genera valores de TC solo dentro de $G$ y copia el resto del volumen de origen. La Figura~\ref{fig:pipeline} resume ese flujo.

\begin{figure}[htbp]
\centering
\small
\fbox{\parbox{0.8\linewidth}{\centering\itshape Compuerta de representación (Objetivo~1), previa a la cadena: su resultado negativo descartó la ruta latente\\ y fijó el dominio de imagen para el sintetizador}}\\[8pt]
\fbox{\parbox{0.8\linewidth}{\centering TC receptora sin osteosíntesis (subconjunto local de CTPelvic1K)}}\\[2pt]
$\downarrow$\\[2pt]
\fbox{\parbox{0.8\linewidth}{\centering Máscaras anatómicas (TotalSegmentator) y medición del corredor óseo\\ en el marco de referencia de Kaiser et al.}}\\[2pt]
$\downarrow$\\[2pt]
\fbox{\parbox{0.8\linewidth}{\centering \textbf{Muestreador de colocación} (Objetivo~2): poses perturbadas alrededor del eje del corredor,\\ evaluadas con SAP frente a las dos distribuciones de la referencia clínica}}\\[2pt]
$\downarrow$\\[2pt]
\fbox{\parbox{0.8\linewidth}{\centering Tornillo paramétrico en la pose muestreada: máscara $M$ y región $G = M \cup B_{\delta}$}}\\[2pt]
$\downarrow$\\[2pt]
\fbox{\parbox{0.8\linewidth}{\centering \textbf{Sintetizador} (Objetivo~3): \emph{inpainting} por difusión en el dominio de imagen dentro de $G$;\\ fuera de $G$ se copia el volumen de origen}}\\[2pt]
$\downarrow$\\[2pt]
\fbox{\parbox{0.8\linewidth}{\centering Evaluación de apariencia con las métricas de Peters et al., frente a la inserción por copia y pegado\\ y frente al protocolo físico}}
\caption{Cadena propuesta. Las cajas en negrita son los dos componentes que se diseñan y se evalúan por separado; la caja superior es la compuerta que fijó el dominio del sintetizador.}
\label{fig:pipeline}
\end{figure}

Lo que se evalúa es la coherencia física y quirúrgica de lo sintetizado, no su utilidad para entrenar redes de segmentación. Quedan fuera de alcance las ablaciones restricción por restricción y la evaluación de segmentación posterior (Dice, HD95). Esta última se excluye por dos condiciones de los datos. La primera es la anotación: de los 75 volúmenes del subconjunto CLINIC-metal de CTPelvic1K, la publicación anota 14 \cite{liu2021ctpelvic1k}, y no está verificado que esas máscaras estén disponibles localmente. La segunda es la cobertura: localmente se dispone de 178 de los 1\,184 volúmenes publicados (Sección~\ref{sec:datos}). Un resultado de segmentación descansaría, por tanto, en una cohorte pequeña y parcialmente anotada.

El estado de ejecución de cada objetivo condiciona lo que este capítulo puede describir como hecho. La compuerta del Objetivo~1 se ejecutó y dio un veredicto negativo; el muestreador y SAP del Objetivo~2 se ejecutaron sobre las cohortes primaria y de sensibilidad. El sintetizador del Objetivo~3 tiene formulación y conjunto de entrenamiento definidos \GAPDATO{la cadena completa, de la pose sobre el corredor medido y el rasterizado de $M$ a la generación de la apariencia sobre una pelvis sin osteosíntesis, se ejecutó sobre dos pacientes de validación; no existe aún resultado del Objetivo~3 porque el modelo final no está elegido (Sección~\ref{sec:sintetizador})}. El Objetivo~3 se rediseñó después de conocer el veredicto del Objetivo~1; la regla de la compuerta se fijó antes de correr la prueba que decide y no se modificó.

\section{Datos y auditoría de la cohorte}\label{sec:datos}

La fuente de anatomía y de observaciones reales de artefacto es CTPelvic1K, que declara 1\,184 volúmenes de TC pélvica \cite{liu2021ctpelvic1k}. Localmente se dispone de 178 de ellos, repartidos en dos carpetas: \texttt{dataset6}, con 103 volúmenes, y \texttt{dataset7}, con 75. La correspondencia de esas carpetas con los subconjuntos CLINIC y CLINIC-metal no está declarada en los encabezados de los archivos: se infiere por coincidencia de conteo, espaciado y dimensiones con la tabla de la publicación. Los recuentos de este trabajo describen el inventario local, no la colección completa.

El artículo del conjunto no reporta los parámetros de adquisición ni de reconstrucción, incluida una posible reducción de artefactos metálicos (MAR, \emph{metal artifact reduction}) del fabricante o una reconstrucción monoenergética. Selles et al.~\cite{selles2024marreview} describen que ambas cambian la apariencia del artefacto. En consecuencia, las observaciones reales de artefacto, cuya reconstrucción es desconocida, se usan para comparar la apariencia y no como verdad física.

El nombre de un subconjunto no establece el tipo de implante ni la disponibilidad de anotación, de modo que la cohorte se construyó por auditoría. Un cribado por HU prioriza la revisión: se marca todo volumen con al menos un vóxel sobre 2500~HU, sin filtro de tamaño. Ese umbral es una heurística de cribado, no un criterio validado de identificación, y no forma parte de las métricas de Peters et al. La autora revisó en vistas tridimensionales los 178 volúmenes y anotó tipo y cantidad de material; la elegibilidad para entrenamiento exige además revisión multiplanar de cortes \GAPDATO{constancia de que la revisión multiplanar completa de cortes está hecha; el registro de datos solo acredita la revisión tridimensional}.

La unidad de independencia es el paciente, no el volumen. Los duplicados exactos se detectan por la huella criptográfica (SHA256) del volumen completo y los parciales por la huella de cada corte, porque la primera no ve subconjuntos ni campos de visión solapados del mismo estudio. Esas dos huellas, más una observación visual que unió dos adquisiciones del mismo implante, agrupan los 178 volúmenes en 168 pacientes. Los 168 pacientes se reparten en tres grupos: 65 con material ortopédico (grupo~1), 37 con objetos no ortopédicos (grupo~2) y 66 sin objeto metálico (grupo~3). Los objetos no ortopédicos son, por ejemplo, material extracorpóreo o dispositivos intrauterinos. La partición trabaja sobre 179 unidades de volumen, porque dos campos de visión solapados del mismo estudio se unen sin interpolar en una unidad adicional. De esas unidades, 10 quedan fuera de uso y una se reserva como par de reproducibilidad del mismo paciente.

La \textbf{regla de aislamiento estricto} prohíbe que un paciente contribuya a la vez al entrenamiento y a la evaluación. En el Objetivo~1, la partición se hizo por paciente y estratificada por subconjunto y por presencia de metal. Separa tres conjuntos: uno de entrenamiento, uno de validación con 8 casos y uno de prueba con 34 pacientes. El Objetivo~3 reutiliza esa partición: sus pares de entrenamiento se forman solo con volúmenes de pacientes de entrenamiento con metal. El Objetivo~2 no entrena ningún modelo, por lo que su cohorte no se divide; se construye con pacientes sin osteosíntesis (grupos~2 y~3), porque el corredor debe medirse sobre hueso sin implante.

La cohorte del Objetivo~2 resulta de un embudo de dos pasos. De 103 pacientes sin osteosíntesis, 91 tenían localizado el nivel S1 y 72 superaron el control de nivel; esos 72 forman la cohorte primaria. El control de nivel excluye un caso en dos situaciones. La primera es que el nivel S1 de las máscaras anatómicas (Sección~\ref{sec:corredor}) no coincida con el que da la localización automática de referencias anatómicas (Sección~\ref{sec:marco-metal}). Esa comparación enfrenta el techo de la máscara de S1, etiqueta de vértebra completa que incluye el arco posterior, con el punto del platillo superior que da la localización. La segunda es que la máscara de S1 toque el borde del campo de visión o esté vacía. La discordancia de nivel excluyó a 11 de 34 pacientes con objetos no ortopédicos y a 7 de 57 pacientes sin objeto, y un paciente más se excluyó porque S1 tocaba el borde del campo de visión. La cohorte de sensibilidad son los 49 pacientes sin objeto que superan el mismo control.

El criterio de selección fue la ausencia de osteosíntesis y no la ausencia de fractura, de modo que la fractura se cribó después. Un cribado ciego de fractura sobre 30 casos, 15 de corredor estrecho y 15 de control, halló fractura confirmada en 20 de los 30 (67\,\%), más un caso dudoso que se reporta aparte y no se cuenta como fractura. La proporción es idéntica en los dos grupos: 10 de 15 corredores estrechos y 10 de 15 controles. El cribado lo hizo un médico recién egresado y sin especialidad, sin conocer a qué grupo pertenecía cada caso; es un cribado válido y no una lectura de especialista, y así se usa en este trabajo. Juzgó 18 casos sobre láminas fijas y 12 sobre el volumen completo, que abrió en un visor cuando la lámina no bastaba. Ese paso al volumen lo motivó una sospecha, de modo que la detección pudo ser más sensible donde ya se sospechaba algo, pero se repartió de forma pareja entre los grupos: 7 corredores estrechos y 5 controles. La cohorte receptora no es, por tanto, una colección de pelvis sanas. Ninguna cifra del corredor depende de ello, porque el corredor se mide sobre cada volumen tal como está, pero sí cambia cómo se compara esa anatomía con la de la referencia clínica (Sección~\ref{sec:amenazas}).

Los 65 pacientes con material ortopédico cumplen cuatro funciones y ninguna de ellas es la de aportar geometría de implante. Sirven de contraste descriptivo para las colocaciones, de comparación de apariencia para la síntesis, de fuente de pares de entrenamiento del sintetizador y, en el único paciente adquirido dos veces con el mismo implante, de par de reproducibilidad. La geometría del implante es paramétrica (Sección~\ref{sec:geometria}), y la auditoría local motiva esa elección. Los tornillos extraídos con un umbral fijo de HU aparecen fragmentados, con diámetros de fuste por debajo de los calibres publicados, y su geometría cambia entre dos adquisiciones del mismo implante.

\section{Evaluación de la representación multiventana}\label{sec:obj1}

La compuerta del Objetivo~1 existe porque un modelo de difusión que opera en el espacio latente de un autoencoder solo puede generar lo que ese autoencoder es capaz de representar. Si el paso de ida y vuelta (HU $\rightarrow$ representación $\rightarrow$ HU) pierde los valores del hueso denso, los HU del hueso quedan sesgados en el volumen sintetizado aunque el generador no tenga error. La compuerta se definió, por eso, como obligatoria, de tipo \emph{Go/No-Go}, y previa al Objetivo~3. Su regla operativa se fijó por escrito antes de correr la prueba que decide la compuerta. Una exploración previa del autoencoder preentrenado sobre los 178 volúmenes locales, incluidos los de los pacientes que después formaron la partición de prueba, ya había quedado por encima del criterio en todos ellos. La regla se fijó, por tanto, conociendo un resultado exploratorio desfavorable.

La cadena evaluada convierte cada corte en una \textbf{codificación multiventana} de tres canales, cada uno con una ventana de HU distinta. Esa codificación pasa por el codificador y el decodificador del autoencoder variacional de Stable Diffusion 1.5, y los HU se recuperan sin usar la imagen original. La lectura toma, en cada vóxel, el canal más estrecho que no esté saturado. El error se mide dentro del hueso, delimitado por el mismo umbral de 150~HU que usa el protocolo de evaluación adoptado \cite{peters2025hybrid,haneda2025aapm}. Para un paciente $i$ con conjunto de vóxeles óseos $\Omega_i$, el error es el error absoluto medio (MAE, \emph{mean absolute error}):
\begin{equation}
\mathrm{MAE}_i = \frac{1}{|\Omega_i|} \sum_{v \in \Omega_i} \left| x_v - \hat{x}_v \right|,
\label{eq:mae}
\end{equation}
donde $x_v$ es el valor original en HU y $\hat{x}_v$ el valor recuperado tras la ida y vuelta.

El umbral de 25~HU no proviene de un criterio publicado, porque no se encontró ningún umbral de aprobación sobre exactitud de los valores de TC. Se ancla en los errores que reportan tres métodos de MAR en la misma escala. Karageorgos et al.~\cite{karageorgos2024ddpm} reportan una raíz del error cuadrático medio (RMSE, \emph{root mean square error}) de 20.2~HU para el algoritmo que ancla la escala del protocolo de Peters et al.~\cite{peters2025hybrid}. Para su modelo de difusión en el dominio del sinograma, evaluado en el dominio de imagen, reportan 12.3~HU, y Yun et al.~\cite{yun2026simulationdriven} reportan 12.74~HU para un modelo de difusión latente. Como para los mismos residuos el RMSE nunca es menor que el MAE, la comparación fija un orden de magnitud y no una equivalencia entre las dos métricas.

El diseño cruza dos variantes del autoencoder con tres codificaciones, lo que da seis combinaciones. Las variantes son el autoencoder preentrenado y una versión con el decodificador adaptado y el codificador congelado, con un decodificador adaptado por cada codificación. Cada decodificador se ajustó con cortes codificados de los pacientes de entrenamiento, separados de los de prueba \GAPDATO{parámetros del ajuste del decodificador (número de cortes, función de pérdida, optimizador y criterio de parada): la pérdida L1 y el optimizador AdamW constan en el código del experimento, y el resto solo como valores por omisión de ese código, no en un registro}. La primera codificación usa las ventanas reportadas por Wang et al.~\cite{wang2025adaptiveweighting}: $[-1000, 2000]$, $[-320, 480]$ y $[-160, 240]$~HU. La segunda eleva el techo de la ventana ancha a 20\,000~HU, y la tercera sustituye la ventana ancha por una compresión arcoseno hiperbólico sobre $[-1000, 20\,000]$~HU.

Una combinación aprueba cuando la media, sobre los 34 pacientes de prueba, del $\mathrm{MAE}_i$ de la Ecuación~\eqref{eq:mae} es menor que 25~HU, y la compuerta aprueba si lo hace alguna de las seis. Esa multiplicidad no se corrige; en su lugar, cada veredicto se acompaña del intervalo de confianza (IC) del 95\,\% por remuestreo de pacientes, y un aprobado cuyo límite superior alcance 25~HU se reporta como marginal. Si aprueban varias combinaciones, la que pasa al Objetivo~3 se elige por un orden fijado a priori y no por su error. Ese orden es arcoseno hiperbólico, luego el techo de 20\,000~HU y por último las ventanas publicadas, cuyo techo de 2000~HU satura el metal; dentro de cada codificación, el decodificador preentrenado precede al adaptado. Si ninguna aprueba, la regla registrada establece que el Objetivo~3 no se ejecuta y que el veredicto negativo se reporta como resultado.

En los volúmenes con metal, el mismo error se reporta además dentro de $B_{\delta}$, de forma descriptiva y fuera de la regla de la compuerta. En esta medición, $B_{\delta}$ son los vóxeles a 12~mm o menos, en distancia euclídea tridimensional, del metal real delimitado por el umbral de 2500~HU, excluido el propio metal. Esa banda contiene las rayas (\emph{streaking}) de baja amplitud que el sintetizador debe generar, y un error restringido al hueso no muestra si el autoencoder las conserva. Rombach et al.~\cite{rombach2022latentdiffusion} diseñan la primera etapa de la difusión latente para eliminar detalle de alta frecuencia, de modo que esa conservación no puede suponerse.

Después del veredicto negativo, la compuerta se extendió a un latente entrenado con TC, el del generador de volúmenes de TC de Guo et al.~\cite{guo2025maisi}, con la misma regla y los mismos pacientes de prueba. Antes de correr ese modelo se comprobó su rango de salida, porque saturar las intensidades en ese rango fija una cota inferior del error que ningún peso puede mejorar. Esa cota se calcula saturando los volúmenes de prueba al rango del modelo, sin modelo alguno, y si ya supera el criterio la evaluación no se corre. El veredicto de la compuerta y el de esta extensión se presentan en el capítulo de resultados.

\section{Muestreador de colocación quirúrgicamente restringido}\label{sec:muestreador}

El muestreador no busca una trayectoria óptima: perturba una trayectoria ya medida sobre la anatomía de cada paciente. La razón es que no existe una pose canónica que reproducir. Zwingmann et al.~\cite{zwingmann2009navigated} reportan distribuciones de malposición distintas para una serie operada con navegación y otra con técnica convencional; esas dos series forman la \textbf{referencia clínica} de este trabajo. Además, la orientación de las superficies óseas que limitan el corredor tuvo, entre especímenes cadavéricos, coeficientes de variación de 7 a 25\,\% en la mayoría de las superficies y de hasta 140\,\% en el ala superior de S1 \cite{ziran2007fluoroscopic}. El muestreador produce entonces una distribución de poses preinscrita (Sección~\ref{sec:poses}). Cada pose recibe un \textbf{grado de brecha cortical}, que mide cuánto sobresale el implante del hueso (Sección~\ref{sec:sap}), y la distribución de grados que resulta se compara con la referencia clínica sin haberse ajustado a ella.

\subsection{Marco de referencia y medición del corredor}\label{sec:corredor}

Las poses se expresan en el marco de referencia de Kaiser et al.~\cite{kaiser2014dysmorphism}, que es computable sobre un volumen de TC. El volumen se reformatea a lo largo del eje sacro, perpendicular al platillo superior de S1. La angulación del corredor se mide en el plano coronal contra la línea que une las crestas ilíacas, y en el plano axial contra la línea que une las espinas ilíacas posteriores. El marco incluye además una regla de longitud útil que exige al menos 5~mm hasta la cortical a cada lado del eje. Este trabajo lee ese \textbf{margen cortical de longitud útil}, $h$, como una distancia radial medida en el plano perpendicular al eje, y no como una holgura añadida al tornillo. Con esa lectura, $h$ coincide con el radio del corredor de 10~mm de Kaiser et al.

La geometría del corredor se mide sobre máscaras anatómicas y no sobre un umbral de HU. A lo largo del eje del corredor más ancho, la mediana de la fracción de vóxeles a 150~HU o menos dentro de las máscaras sacras fue 0.42, con el mismo umbral de hueso del protocolo adoptado (Sección~\ref{sec:obj1}). Una definición de hueso por HU dejaría fuera, en la mediana, esa fracción del corredor.

Las máscaras provienen de TotalSegmentator 2.18.0 \cite{wasserthal2023}, construido sobre nnU-Net \cite{isensee2021}, con la tarea \texttt{total}. Antes de segmentar, la herramienta delimita la región de interés con un modelo de recorte de menor resolución. El \textbf{recorte por defecto} usa un modelo de 6~mm; el \textbf{recorte alternativo}, activado por la opción \texttt{-{}-robust\_crop}, uno de 3~mm. El primero es el principal y el segundo se usa como análisis de sensibilidad. Cada caso pasa el control de nivel descrito en la Sección~\ref{sec:datos}, y se eliminan los componentes conexos menores que el 0.1\,\% de su estructura. La publicación citada evalúa un modelo de 104 estructuras y no reporta exactitud para la etiqueta de S1, por lo que esa exactitud no se asume para la versión y las clases usadas aquí.

La \textbf{envolvente ósea} es la unión de las máscaras de sacro, vértebra S1 y ambos coxales, con un cierre morfológico de 2~mm y sin relleno de cavidades cerradas. La decisión que la definía con ese relleno se corrigió, porque las máscaras de TotalSegmentator son volúmenes sólidos y no dejan hueco que tapar. Aplicar el relleno no cambia el diámetro del corredor en ninguno de los 72 pacientes de la cohorte primaria. El diámetro del corredor se define como el doble de la distancia libre mínima a esa envolvente a lo largo del eje, excluidos 8~mm en cada extremo (Sección~\ref{sec:sap}). Por cada caso, la búsqueda devuelve el eje del corredor más ancho, con su centro, su dirección y su longitud \GAPDATO{descripción completa del algoritmo de búsqueda del eje del corredor; hoy consta en el código y, en parte, en el registro de implicancias del proyecto}. Si en el volumen hay un implante, este se trata como espacio ocupado mediante un umbral de semimáximo local por objeto, y un umbral fijo de 2500~HU da la cota superior del corredor ocupado. Esos volúmenes se reportan como contraste aparte y no se agregan a la cohorte del corredor.

Un corredor es viable cuando su diámetro $D$ admite el calibre del tornillo más una \textbf{holgura radial} a cada lado, $D \geq d + 2\epsilon$, donde $d$ es el calibre nominal del tornillo paramétrico (Tabla~\ref{tab:geometrias}) y $\epsilon$ vale entre 1 y 2~mm. La holgura radial se toma de Kaiser et al.~\cite{kaiser2014dysmorphism}, que justifican su corredor de 10~mm como holgura de 1 a 2~mm alrededor de un tornillo de 6.3 a 8~mm. Leerla como holgura radial por lado es una operacionalización propia, porque la publicación no indica si la holgura se cuenta por lado o en total. Kaiser et al. eligen ese corredor de 10~mm como un valor conservador y lo atribuyen a trabajos previos. McLaren et al.~\cite{mclaren2021corridor} lo toman de Kaiser et al. y lo emplean con un procedimiento reproducible de medición del diámetro, pero no lo validan como umbral de seguridad. El umbral fijo de 10~mm se conserva solo como convención de comparación con esa literatura y no como estándar clínico heredado.

Con el recorte por defecto, el diámetro del corredor tuvo en la cohorte primaria una mediana de 9.5~mm (rango intercuartílico de 7.4 a 11.7~mm). Con ese mismo recorte, el criterio $D \geq d + 2\epsilon$ se cumplió en el 65.3\,\% de los 72 pacientes en su extremo menos exigente ($d = 6.5$~mm, $\epsilon = 1$~mm). En el extremo más exigente ($d = 8.0$~mm, $\epsilon = 2$~mm, Tabla~\ref{tab:geometrias}) se cumplió en el 23.6\,\%. Con el recorte alternativo, esos porcentajes fueron 62.5 y 20.8\,\%. La convención de 10~mm se cumplió en 29 de 72 pacientes con el recorte por defecto y en 27 con el alternativo. La diferencia entre recortes es incertidumbre de la medición, no variabilidad anatómica. La limpieza de componentes no cambió el diámetro ni su viabilidad en ninguna fracción probada hasta el 5\,\%; la fracción de 0.1\,\% que se usa, como la regla de ocupación, se fijó después de inspeccionar estos resultados.

Las cifras principales del Objetivo~2 usan el recorte por defecto, y la decisión que lo fijó quedó condicionada a revisar los 16 casos marcados por diferencia de diámetro entre recortes o por desplazamiento de la caja de una estructura. Esa revisión se hizo sobre láminas, por la autora con apoyo de un médico egresado y sin especialidad, y cubrió los 16 casos. De ellos, 14 salieron conformes, 2 fallaron con los dos recortes y ninguno falló solo con el de 6~mm, que es el único veredicto que habría reabierto la decisión. El recorte por defecto se mantiene, ahora con la condición evaluada y no solo declarada; los 2 casos que fallan con ambos recortes son errores de segmentación y no de recorte \GAPDEC{tratamiento de los dos casos con error de segmentación anotado en la revisión de láminas}.

\subsection{Aplicabilidad del marco en volúmenes con metal}\label{sec:marco-metal}

El marco de Kaiser et al. se apoya en referencias anatómicas de la unión lumbosacra, que es donde se concentran las rayas del artefacto metálico. Ninguna de las fuentes revisadas dice si esas referencias siguen siendo localizables con metal, de modo que se cuantificó en la cohorte local. Las referencias se localizaron con una heurística automática \GAPDATO{descripción de la heurística de localización de las referencias anatómicas, que hoy consta en el código y, en parte, en las notas del experimento}, y una región se contó como contaminada cuando su fracción de rayas oscuras o de metal superaba el máximo observado en 69 pacientes de \texttt{dataset6} que la revisión tridimensional clasificó sin objeto. Ese conjunto de calibración no coincide con el grupo~3 de la Sección~\ref{sec:datos}, que tras el reparto por grupos reúne 66 pacientes. Un revisor clínico, cirujano otorrinolaringólogo y ciego a la auditoría automática, revisó el nivel S1 en cada paciente.

De 65 pacientes con material ortopédico, 57 tenían las cinco referencias localizadas y dentro del campo de visión. De los 8 restantes, 7 tenían las crestas ilíacas truncadas por el campo de visión. El marco fue computable con el nivel S1 correcto según el revisor clínico en 48 pacientes, y quedó libre de contaminación en las cinco regiones en 29. En los 17 pacientes en que el censo morfológico del metal identificó un candidato a tornillo transilíaco-transsacro (Sección~\ref{sec:geometria}), el nivel S1 fue correcto en todos y su región estaba contaminada en 3. La localización heurística de crestas y espinas ilíacas no se ha revisado clínicamente.

La autora volvió a juzgar el nivel S1 sobre las mismas 61 láminas sagitales. El acuerdo con el revisor clínico fue de 61 de 61 casos, con dos límites: el segundo lector es la autora y no un segundo clínico, y solo los juicios de nivel son independientes. Las notas de texto libre se copiaron una vez comprobada la coincidencia, de modo que no aportan evidencia independiente. Las revisiones de láminas registran, además, qué estructura se acepta como S1 cuando la máscara se extiende por detrás del cuerpo vertebral. En 12 de los 16 casos revisados por el recorte (Sección~\ref{sec:corredor}) es la cresta sacra media, que la etiqueta \texttt{vertebrae\_S1} de TotalSegmentator contiene por diseño, porque etiqueta la vértebra completa y no solo su cuerpo.

\subsection{Geometría del implante}\label{sec:geometria}

Los implantes se modelan como tornillos rígidos paramétricos en lugar de extraerse de los volúmenes con metal. La trayectoria que este trabajo mide y muestrea va de la cortical externa de un ilion, cruzando las dos articulaciones sacroilíacas, a la cortical externa del ilion contralateral. El implante que representa es, por tanto, un tornillo transilíaco-transsacro, y le aplican las tolerancias angulares de McLaren et al.~\cite{mclaren2021corridor}, medidas sobre ese mismo corredor. Se mantienen separadas tres geometrías, una para decidir la viabilidad, otra para sintetizar y otra para medir la brecha, que la Tabla~\ref{tab:geometrias} resume. La cifra nominal de estos tornillos es el diámetro de la rosca, no el del cuerpo. Un cilindro uniforme de 6.5 a 8.0~mm tendría, a lo largo del corredor, más sección de metal que el cuerpo de 4.91~mm de la Tabla~\ref{tab:geometrias}. El calibre nominal solo se usa para la viabilidad, y la máscara que se sintetiza es el cuerpo del tornillo.

\begin{table}[htbp]
\centering
\small
\caption{Las tres geometrías del implante y su propósito.}
\label{tab:geometrias}
\begin{tabular}{p{0.2\linewidth} p{0.2\linewidth} p{0.5\linewidth}}
\hline
\textbf{Geometría} & \textbf{Propósito} & \textbf{Valor y respaldo} \\
\hline
Calibre nominal & Viabilidad del corredor & 6.5 a 8.0~mm \cite{gardner2010safezones}; 6.3 a 8~mm \cite{kaiser2014dysmorphism} \\
Máscara del implante $M$ & Lo que se sintetiza (Objetivo~3) & Cilindro uniforme de 4.91~mm, mediana del diámetro del cuerpo medida en la cohorte; fuste de 4.8~mm en la guía del fabricante \cite{synthes2003guide}; 4.9~mm de fuste y 4.7~mm de núcleo bajo la rosca \cite{gardner2015screw} \\
Calibre de medición & Lo que mide la brecha en SAP & 7.0~mm, el de los tornillos canulados de la referencia clínica \cite{zwingmann2009navigated}; 4.91 y 7.3~mm como sensibilidad \\
\hline
\end{tabular}
\end{table}

El cuerpo de 4.91~mm es la mediana medida en la cohorte local sobre 79 componentes metálicos alargados y finos de 43 volúmenes. Los seleccionó el \textbf{censo morfológico del metal}, un filtro geométrico y no clínico: vóxeles sobre 2500~HU, longitud de al menos 30~mm y anchos de 12~mm o menos. Ese censo no afirma que los componentes sean tornillos transilíaco-transsacros, sino que son candidatos por su forma. La cifra está próxima al fuste de 4.8~mm que imprimen la guía técnica del fabricante \cite{synthes2003guide} y el modelo de elementos finitos de Zhu y Liao~\cite{zhu2022optimalposition}. La guía da ese mismo fuste para los tornillos de 6.5 y 7.3~mm, y el modelo usa 7.3~mm solo en una rosca de 16~mm.

La cabeza, cuando se modela, mide 8.0~mm de diámetro y 4.5~mm de altura. El diámetro lo reportan Sayres et al.~\cite{sayres2014comparison} y lo imprime también el catálogo de otro fabricante \cite{doublemedical2021trauma}; la altura, entre las fuentes revisadas, solo la dan Sayres et al. Que la cabeza quede sin avellanar es un supuesto propio; la cabeza y la rosca entran como variantes de sensibilidad.

El calibre de medición lo fija la referencia clínica y no se elige. Zwingmann et al.~\cite{zwingmann2009navigated} asignaron sus cuatro grados sobre tornillos canulados de 7.0~mm, y la profundidad de perforación crece con el radio. Medir las colocaciones simuladas con otro calibre introduciría en la distancia de comparación un sesgo sistemático que no se podría separar del error de colocación. La longitud del tornillo tampoco procede de un rango publicado: se acota en cada volumen por el corredor medido, con la regla de longitud útil de Kaiser et al.

Tres parámetros quedan sin fijar por la documentación del propio fabricante y se declaran en lugar de suponerse. El espesor de la arandela se toma como 1.5~mm del catálogo de otro fabricante \cite{doublemedical2021trauma}. Las dos fuentes de fabricante \cite{synthes2003guide,doublemedical2021trauma} son documentación técnica no revisada por pares. El diámetro de la canulación queda como parámetro libre y se evalúa el tornillo macizo frente al hueco, porque el valor que circula en listados de distribuidores solo aparece en la guía del fabricante como canulación de brocas y destornilladores. La aleación tampoco se codifica: el mismo sistema se fabrica en acero inoxidable 316L y en Ti-6Al-7Nb \cite{synthes2003guide}. Una máscara binaria no distingue entre ambos, de modo que el sintetizador se entrena con la apariencia de una mezcla de materiales.

\subsection{Distribución de poses preinscrita}\label{sec:poses}

La distribución de poses y la métrica se congelaron por escrito el 22 de septiembre de 2026, antes de calcular cualquier distancia contra la referencia clínica. Esa \textbf{preinscripción} (\emph{pre-registration}) impide que la comparación sea un ajuste: si la distribución se hubiera modificado después de ver la distancia, se compararía el modelo contra los mismos datos usados para construirlo. Cualquier cambio posterior se registra como desviación, con fecha y motivo.

Para cada caso, la pose base es el eje del corredor medido en ese volumen, anclado en el punto medio de su tramo de corredor, y de ese mismo tramo sale la longitud del implante. El anclaje corrige un desajuste de la búsqueda: el punto desde el que se lanza la búsqueda no siempre coincide con el centro del tramo óseo encontrado. Sin él, un implante simétrico alrededor de ese punto quedaría descentrado, y la identidad ``implante sobre el eje del corredor, brecha nula'' dejaría de cumplirse.

Sea $\mathbf{c}$ el centro anclado y $\mathbf{u}$ la dirección unitaria del eje. Cada pose aplica dos perturbaciones en el plano perpendicular a $\mathbf{u}$, ambas con dirección uniforme en ese plano y magnitud seminormal truncada en tres desviaciones estándar:
\begin{equation}
\mathbf{c}' = \mathbf{c} + \rho\,\mathbf{e}_{\phi}, \qquad \rho = |N(0, \sigma_t)|, \qquad \mathbf{u}' = R_{\psi}(\alpha)\,\mathbf{u}, \qquad \alpha = |N(0, \sigma_a)|,
\label{eq:pose}
\end{equation}
donde $\mathbf{e}_{\phi}$ es un vector unitario perpendicular a $\mathbf{u}$ con ángulo $\phi$ uniforme, y $R_{\psi}(\alpha)$ es una rotación de ángulo $\alpha$ alrededor de un eje perpendicular a $\mathbf{u}$ de orientación $\psi$ uniforme. La Tabla~\ref{tab:preinscripcion} recoge los parámetros.

\begin{table}[htbp]
\centering
\small
\caption{Parámetros preinscritos de la distribución de poses.}
\label{tab:preinscripcion}
\begin{tabular}{p{0.3\linewidth} p{0.22\linewidth} p{0.38\linewidth}}
\hline
\textbf{Parámetro} & \textbf{Valor} & \textbf{Procedencia} \\
\hline
Margen cortical $h$ & 5.0~mm & Regla de longitud útil de Kaiser et al.~\cite{kaiser2014dysmorphism} \\
$h = 2\sigma$ & convención & Propia y declarada; no es un dato medido \\
$\sigma_t$ & $h/2 = 2.5$~mm & Derivado \\
$\sigma_a$ & $\frac{1}{2}\arctan\!\left(\frac{h}{L/2}\right)$, por caso & Derivado; $2.07^{\circ}$ con la mediana de longitud del corredor, $L = 138$~mm \\
Truncamiento & $3\sigma$ & Declarado; acota la magnitud de cada perturbación \\
Poses por caso & 50 & Declarado \\
Semilla & 20260922, derivada por caso & Fija; las poses no se vuelven a sortear \\
\hline
\end{tabular}
\end{table}

Toda la escala deriva de una sola constante independiente, el margen cortical de 5~mm de la regla de longitud útil (Sección~\ref{sec:corredor}). La convención propia fija ese margen en dos desviaciones estándar: en distancia para $\sigma_t$ y en el ángulo que produce para $\sigma_a$ \GAPDEC{justificación de la convención $h = 2\sigma$, que fija la escala de todas las perturbaciones y no consta en el repositorio}. El margen $h$ es una distancia radial al eje, es decir, se mide en el mismo plano perpendicular en que actúan las perturbaciones. La dispersión angular $\sigma_a$ es ese mismo margen expresado como el ángulo que lo produce en la punta del corredor de cada caso. Ningún parámetro se toma de Zwingmann et al.~\cite{zwingmann2009navigated}, incluida la tolerancia angular de $4^{\circ}$ que su introducción cita de una tercera fuente: aparece en el mismo artículo contra el que se puntúa el resultado. La perturbación es isótropa porque privilegiar una dirección anatómica de error exigiría un dato clínico de direccionalidad que no se tiene. Fijarla a juicio añadiría un parámetro sin dato que lo sostenga.

La semilla se deriva del identificador de cada caso. La cohorte de sensibilidad es un subconjunto de la primaria, y con un generador que avanzara de caso en caso el mismo paciente recibiría poses distintas en las dos corridas. Con la semilla por caso, cada paciente recibe las mismas poses en cualquier cohorte que lo contenga, y las dos cohortes difieren solo en su composición.

Antes de la corrida se fijaron tres comprobaciones que la implementación debía cumplir. Son controles de consistencia de la implementación y no condiciones de fallo del muestreador como modelo; las dos primeras forman parte de los seis controles de SAP (Sección~\ref{sec:sap}), y la tercera se comprueba en cada corrida. La pose sin perturbar debe dar brecha nula con el diámetro de su propio corredor, porque es el eje del corredor. La brecha no debe decrecer al ensanchar el cilindro. Ninguna pose puede quedar sin calificar, porque descartar las que no atraviesan hueso eliminaría del denominador las peores colocaciones y bajaría la distancia reportada sin que la colocación mejorara.

Bajo el corredor de S1, la búsqueda devuelve un segundo corredor óseo cuyo nivel vertebral no es constante. En una submuestra de 18 volúmenes estratificada por profundidad, una segunda lectura de la autora, ciega a la profundidad medida, lo ubicó en S2 en 13 volúmenes, en S3 en 4 y en S4 en 1. La profundidad no determina el nivel: un corredor a 39~mm bajo S1 se juzgó S2, mientras otros a 36 y 33~mm se juzgaron S3. De ahí que se nombre ``segundo corredor bajo S1'' y no S2. Sus resultados son descriptivos, porque ninguna de las fuentes revisadas reporta una distribución ordinal clínica por debajo de S1, y la preinscripción cubre solo S1 \GAPDATO{preinscripción y corrida del muestreo de poses en el segundo corredor bajo S1; su eje ya está medido en 69 de los 72 pacientes de la cohorte primaria}.

La densidad ósea entra como componente reportado de SAP y no como condicionante del muestreo. Si condicionara el muestreo, se volvería un parámetro ajustable más y reabriría la circularidad que la preinscripción cierra. La motivación es la heterogeneidad de la masa ósea pélvica que describen Arand et al.~\cite{arand2019pelvicring} en su modelo estadístico poblacional, por ejemplo los valores de gris bajos en el ala sacra \GAPDEC{fuente y definición operativa de la fracción por zona de densidad: la decisión que la fija como componente de SAP la calcula sobre una sonda de densidad que una decisión anterior retiró como evidencia de densidad, y la definición implementada solo consta en el código}.

Los fenotipos de morfología sacra no se usan como variable de estratificación, porque su efecto sobre el tamaño de la zona segura no es consistente entre estudios. Gardner et al.~\cite{gardner2010safezones} encuentran una zona segura de S1 menor en sacros dismórficos, pero recogen un estudio previo que no halló diferencia entre sacros normales y dismórficos. El muestreador mide, en cambio, la geometría del corredor directamente en cada volumen.

\section{Sintetizador condicionado en el dominio de imagen}\label{sec:sintetizador}

El sintetizador se formula como \emph{inpainting} de la región de generación $G = M \cup B_{\delta}$ mediante un modelo de difusión condicional 2.5D en el dominio de imagen, es decir, sobre valores de TC por vóxel, sin espacio latente ni datos de proyección. El modelo recibe el parche con $G$ borrada, la máscara del implante $M$ y la máscara de la banda, y genera valores de TC solo dentro de $G$. Todo lo que queda fuera de $G$ se copia del volumen de origen, de modo que la preservación fuera de la banda se cumple por construcción y no por medición. Las entradas y salidas usan la codificación multiventana como representación de canales, sin la etapa de compresión del autoencoder. La codificación elegida entre las tres de la compuerta, la arquitectura de la red y los parámetros de entrenamiento y muestreo constan en un diseño que todavía no se congeló \GAPDEC{congelar la preinscripción del diseño del sintetizador: codificación multiventana (el borrador supone la de arcoseno hiperbólico), arquitectura, los demás parámetros de entrenamiento, muestreo, número de semillas por caso y margen de equivalencia}.

La codificación multiventana acota por abajo los valores de TC representables, porque su ventana ancha empieza en $-1000$~HU (Sección~\ref{sec:obj1}). Qué se hace con ese suelo lo decide una regla preinscrita. La regla mide, sin modelo y en la banda $B_{\delta}$ de los pacientes con metal de entrenamiento y validación, qué parte de la diferencia entre las colas del 5\,\% de los HU pierde la ida y vuelta por la representación. Esa pérdida se evalúa en el percentil~5 de cada paciente y no en la mediana, una elección hecha después de ver su distribución. Si es menor que el margen de equivalencia $\Delta$ (Sección~\ref{sec:apariencia}), el suelo se declara como limitación de alcance; si es mayor, se cambia la representación \GAPDATO{desenlace de la regla preinscrita sobre el suelo de $-1000$~HU, que depende del margen de equivalencia $\Delta$, aún no medido}.

Esta formulación sustituye a la inicial, que era difusión latente sobre la base preentrenada Stable Diffusion 1.5 guiada con ControlNet. ControlNet presupone una base preentrenada congelada y su espacio latente \cite{zhang2023controlnet}, y el Objetivo~1 midió que ese latente no conserva los valores de TC del hueso con la tolerancia exigida. Como no se adopta ninguna base de TC preentrenada en el dominio de imagen, la guía por ControlNet pierde su objeto.

Los HU de la salida del sintetizador se leen de su codificación multiventana con una regla propia del Objetivo~3, distinta de la del Objetivo~1. La lectura del Objetivo~1 toma el canal más estrecho que no esté saturado (Sección~\ref{sec:obj1}), y con ella un canal estrecho casi saturado sustituía el valor del canal de ventana ancha por el techo de su propia ventana. El sintetizador parte, en cambio, del canal de ventana ancha y lo corrige con cada canal estrecho mediante una mezcla ponderada, en la que un canal pesa menos cuanto más se acerca a la saturación. El peso depende de la distancia del valor normalizado del canal al borde de su ventana: es nulo por debajo de 0.01, crece de forma lineal y es completo desde 0.05. Si los dos límites coinciden en 0.01, que es el umbral de saturación de la lectura del Objetivo~1, la rampa se reduce a un escalón y la mezcla da esa misma lectura, que queda así como caso particular. El límite de 0.05 se fijó sobre el conjunto de validación, con un criterio registrado antes de conocer el resultado \GAPDATO{resultado citable de la calibración de ese límite sobre el conjunto de validación}. La compuerta del Objetivo~1 conserva su lectura, y sus cifras no cambian.

La banda $B_{\delta}$ extiende la síntesis unos 12~mm más allá del borde del implante, para que las rayas y el endurecimiento del haz puedan manifestarse fuera de la máscara. Su ancho es un parámetro de diseño: acota dónde se genera artefacto y trunca a propósito las rayas lejanas, cuya extensión espacial no está medida en la literatura revisada. Ninguna de las fuentes revisadas da un valor para el ancho. Karageorgos et al.~\cite{karageorgos2024ddpm} sí dan el sentido del error en su propio dominio: estudian la sensibilidad de su método a la máscara metálica, con la que obtienen un RMSE de 7.57~HU cuando la máscara es la verdadera. Dilatarla eleva el error a 11.45 y 11.63~HU para factores de 1.4 y 1.8; contraerla lo eleva a 54.82 y 57.69~HU para factores de 0.7 y 0.5. Ese experimento varía, sin embargo, una máscara de reducción de artefactos en el dominio del sinograma y no una banda de generación en el dominio de imagen. Este trabajo lo lee como apoyo a generar más allá de la máscara del implante, pero no como valor para el ancho de 12~mm.

Los pares de entrenamiento se forman con los volúmenes con metal de los pacientes de entrenamiento, donde el implante real aporta a la vez la máscara y el objetivo de apariencia. La máscara de entrenamiento se obtiene con el umbral de 2500~HU, y el semimáximo local por objeto entra como sensibilidad. La unidad de entrenamiento es cada componente conexo de metal y no el corte axial, para que una región de entrenamiento tenga la misma estructura que una región de síntesis, que contiene un solo implante. En la síntesis, la máscara es el tornillo paramétrico de la Sección~\ref{sec:geometria} colocado en una pose muestreada. El entrenamiento se fija en 30\,000 pasos, una elección de hiperparámetro hecha sobre el conjunto de validación y no fijada de antemano. Dos corridas sitúan el mínimo de la curva de validación en una meseta, y ese valor ocupa su centro. Ningún paciente de prueba intervino en la elección, de modo que la regla de aislamiento estricto (Sección~\ref{sec:datos}) se mantiene.

La inclusión en el conjunto de entrenamiento se rige por un criterio fijado antes de entrenar, con tres reglas. El componente debe estar dentro del cuerpo del paciente, y el caso debe tener material ortopédico declarado en la auditoría. El parche, por último, debe contener metal del componente objetivo. A esos parches se añaden todos los parches disponibles que solo contienen banda, que están en razón de 0.62 respecto de los parches con metal. El tamaño y la forma del componente se reportan y no se filtran, porque se asume que la relación entre máscara y artefacto no depende del tipo de implante. Ese supuesto no se verificó. Lo que sí se midió es el costo de filtrar por forma: en una muestra revisada de 40 componentes, solo 5 eran tornillos, de modo que un filtro de forma habría descartado la mayor parte de los ejemplos.

Entre entrenamiento y síntesis hay tres desplazamientos de dominio. Primero, las máscaras de entrenamiento se obtienen umbralizando metal real, cuyo fuste la auditoría midió por debajo de los calibres publicados, mientras que la máscara de síntesis es un cilindro paramétrico liso. Segundo, fuera de $G$ el contexto trae rayas lejanas en los volúmenes de entrenamiento con metal, pero no en los volúmenes sin metal que reciben la síntesis. Tercero, los parches que solo contienen banda son menos frecuentes en el conjunto de entrenamiento de lo que la geometría del corredor predice para la síntesis: 0.62 frente a 1.40 parches de solo banda por parche con metal. El tercero queda así cuantificado. Para el primero falta la métrica \GAPDEC{métrica y análisis con que se cuantifica la diferencia entre la máscara umbralizada de entrenamiento y el cilindro liso de síntesis}. Para el segundo, el salto de los valores de TC al cruzar el borde de la banda ya se midió una vez, sobre un caso, un corte y un modelo sin converger \GAPDEC{preinscribir el bloque que mide la continuidad de los valores de TC en el borde de la banda y su análisis}.

El \textbf{punto de control} (\emph{checkpoint}) del sintetizador que pasa a la evaluación no se elige por la pérdida de validación, porque esa pérdida no ordena los puntos de control por la fidelidad de los HU que generan. Esa misma razón debilita el argumento de los 30\,000 pasos, tomado del mínimo de esa curva: la cifra se mantiene y se revisará cuando el criterio de selección esté medido. Tampoco se elige por la reconstrucción de implantes reales de validación, porque esa tarea premia memorizar implantes parecidos a los de entrenamiento y no es la de este trabajo, que coloca un tornillo donde no había metal. El criterio se mide, por eso, sobre la tarea de síntesis. Se coloca el tornillo paramétrico en una pelvis de validación sin osteosíntesis, se genera con cada punto de control candidato y se comparan dos estadísticos de la apariencia sintética con los mismos estadísticos medidos en los implantes reales de validación. Ninguno de los dos exige una imagen sin metal del paciente, que en un implante real no existe y que la \emph{streak amplitude} sí necesita (Sección~\ref{sec:apariencia}). Los HU sintéticos se leen con la mezcla de canales descrita arriba.

El primer estadístico es el perfil radial de HU alrededor del metal. Se mide en el plano de cada corte, por cáscaras de 0.5~mm de 0 a 12~mm en anillo completo. En cada cáscara se toman la mediana y el percentil~95 de HU, expresados como elevación sobre la mediana del anillo de 12 a 15~mm del mismo paciente. El percentil~95 acompaña a la mediana porque las rayas ocupan las colas de la distribución de HU, que la mediana no registra. Los valores se agregan por corte, después por paciente y después entre pacientes. El segundo estadístico es el histograma de HU dentro de $M$, resumido por paciente en los percentiles~50, 75 y~95 de los vóxeles sobre 2500~HU. La fracción sintética de $M$ que supera 2500~HU se reporta solo como control y no se compara con la real, porque en los implantes reales la máscara se define con ese mismo umbral y esa fracción no informa nada.

Como referencia se usan los implantes reales de los 3 pacientes de validación (Sección~\ref{sec:apariencia}). Gana el punto de control con más cáscaras dentro de la envolvente real, del mínimo al máximo de esos 3 pacientes, en las dos curvas del perfil. Un empate lo decide la menor distancia a la mediana real de los percentiles~50 y~95 dentro de $M$, y un empate que persista lo decide la autora. Con 3 pacientes de referencia, el cotejo puede descartar un punto de control cuyo artefacto cae fuera de lo observado, pero no prueba que el elegido lo reproduzca. Exige además que esos 3 pacientes lleven tornillos, y su tipo de implante no está verificado; si no lo son, el cotejo se revisa \GAPDEC{verificación del tipo de implante de los 3 pacientes de referencia antes del cotejo; el procedimiento para verificarlo es una propuesta preliminar}. Mientras el cotejo no esté medido, ningún punto de control está elegido.

\section{Protocolo de evaluación}\label{sec:evaluacion}

La evaluación separa la colocación de la apariencia porque cada una tiene su propia referencia. La colocación se juzga contra la referencia clínica; la apariencia, contra la \textbf{inserción por copia y pegado} y contra el protocolo físico de Peters et al. La Tabla~\ref{tab:diseno} resume el diseño experimental de los tres objetivos que se evalúan con datos. El Objetivo~4 no tiene fila propia: su producto es la definición de las métricas que usan las filas~2 y~3. La implementación de SAP se verificó con los controles de la Sección~\ref{sec:sap} \GAPDEC{qué evidencia verifica el Objetivo~4 más allá de esos controles, o si se declara que no se evalúa experimentalmente}.

\subsection{Admisibilidad quirúrgica de la colocación}\label{sec:sap}

SAP tiene tres componentes: el grado de brecha cortical de cada pose, la fracción por zona de densidad y la viabilidad del corredor en el marco de Kaiser et al. Solo la distribución de grados se compara con la referencia clínica. La fracción por zona de densidad y la viabilidad se reportan de forma descriptiva, la segunda por caso y con el criterio $D \geq d + 2\epsilon$ de la Sección~\ref{sec:corredor} \GAPDEC{calibre con que SAP evalúa la viabilidad del corredor: la decisión que define SAP fija el calibre nominal de 6.5 a 8.0~mm, y la corrida del Objetivo~2 usó los calibres de 4.91, 7.0 y 7.3~mm con una holgura de 1~mm}. El grado de brecha cortical usa la escala ordinal de cuatro grados que Smith et al.~\cite{smith2006iliosacral} aplican a la fijación iliosacra. Smith et al. suman además una escala angular a la de perforación \GAPDEC{si SAP incorpora la dimensión angular de la escala de Smith et al. o declara que usa solo la de perforación, y con qué justificación}. La brecha de una pose es la protrusión radial máxima del implante fuera de la envolvente ósea segmentada:
\begin{equation}
b = \max_{t \in T} \; \max\bigl(0,\; r + s(\mathbf{c}' + t\,\mathbf{u}')\bigr),
\label{eq:brecha}
\end{equation}
donde $r$ es el radio del calibre de medición, $s$ es la distancia con signo a la envolvente ósea (negativa dentro del hueso) y $T$ es el tramo evaluado a lo largo del eje. La envolvente es una máscara de hueso y no una segmentación de la cortical, así que su borde aproxima la superficie cortical externa. Los límites de grado se fijan por convención: grado~0 si $b = 0$; grado~1 si $b \in (0, 2)$~mm; grado~2 si $b \in [2, 4]$~mm; y grado~3 si $b > 4$~mm.

El tramo $T$ es el propio implante, de longitud fija igual a la del corredor del caso, centrado en la pose y sin los 8~mm de cada extremo. Esa exclusión de extremos existe porque un tornillo transilíaco-transsacro entra y sale por la cortical del ilion por diseño; sin ella, toda trayectoria válida puntuaría grado~3 por sus propios puntos de entrada y salida. Los 8~mm son los mismos que se excluyen al medir el diámetro del corredor (Sección~\ref{sec:corredor}) \GAPDEC{justificación del valor de 8~mm de la exclusión de extremos, que no consta en el repositorio}. La longitud no depende de dónde esté el hueso, porque si dependiera la métrica premiaría a las poses que se salen del corredor, cuyo tramo se acortaría con ellas. Las poses que no atraviesan hueso reciben grado~3 en lugar de descartarse, para que el denominador coincida con el de la referencia clínica, que calificó todos sus tornillos.

Seis controles verificaron la implementación antes de la corrida, entre ellos la brecha nula sobre el eje y la monotonía en el diámetro. Uno se hizo sobre un fantoma de geometría conocida y cinco sobre el eje del corredor de dos volúmenes reales. La resolución de la métrica en el fantoma es de 0.083~mm.

La distribución de grados de las poses se compara con las dos distribuciones ordinales de Zwingmann et al.~\cite{zwingmann2009navigated}, condicionadas por técnica quirúrgica y en S1, nivel que la serie convencional declara y que en la navegada se asume (Sección~\ref{sec:amenazas}). En la serie navegada, los grados~0 a~3 reúnen el 69, 15, 8 y 8\,\% de los tornillos; en la convencional, el 40, 37, 11.5 y 11.5\,\%. Con los grados tratados como equiespaciados por convención propia, la distancia de Wasserstein-1 entre dos distribuciones $P$ y $Q$ sobre los grados $0$ a $3$ es
\begin{equation}
W_1(P, Q) = \sum_{k=0}^{2} \left| F_P(k) - F_Q(k) \right|,
\label{eq:w1}
\end{equation}
donde $F_P$ y $F_Q$ son las funciones de distribución acumulada. Su unidad es el grado, no el milímetro.

Una distancia grande no se corrige moviendo parámetros del muestreador, que no se ajustó a esas distribuciones. El Objetivo~2 no tiene regla de decisión: ningún valor de la distancia se fijó de antemano como fallo del muestreador \GAPDEC{si se fija una escala de lectura de la distancia de Wasserstein-1, o qué resultado contaría como fallo del muestreador}. La comparación enfrenta una población simulada sobre pelvis receptoras con dos series clínicas pequeñas, de 26 tornillos en 24 pacientes y de 35 tornillos en 32 pacientes \cite{zwingmann2009navigated}; no es una prueba de equivalencia. Los grados se cuentan por tornillo. Los recuentos implican que algunos pacientes recibieron más de uno, aunque el artículo declara un solo tornillo por paciente.

Además del análisis preinscrito, se añade un análisis post hoc estratificado por si el diámetro del corredor admite el calibre de medición de 7.0~mm (Tabla~\ref{tab:geometrias}). En un corredor más estrecho que ese calibre, el propio eje del corredor ya perfora, y el grado~0 es inalcanzable por la anatomía receptora tal como se presenta, no por la colocación. Eso ocurre en 15 de los 72 pacientes de la cohorte primaria. Qué estrecha esos corredores no está establecido. El cribado de fractura halló fractura sacra desplazada en 7 de los 15, un indicio sugerente y no concluyente (Sección~\ref{sec:datos}). Una fracción no cuantificable del estrechamiento puede ser, por tanto, patología y no variabilidad anatómica normal. La estratificación se reporta junto a la cifra preinscrita y no la sustituye. Restringir la cohorte a los corredores que admiten el calibre de 7.0~mm se consideró y se rechazó por dos razones. La primera es que seleccionar pacientes por una variable correlacionada con el resultado subiría la fracción de grado~0 por selección y no por colocación, el mismo sesgo que introduciría ajustar el muestreador a la referencia clínica. La segunda, que Zwingmann et al. no reportan haber excluido pacientes de corredor estrecho, de modo que una cohorte restringida dejaría de ser comparable con la suya.

\subsection{Coherencia de apariencia}\label{sec:apariencia}

La apariencia se evalúa con las métricas del protocolo de Peters et al.~\cite{peters2025hybrid}, con sus nombres publicados: \emph{bone integrity}, \emph{metal integrity} y \emph{streak amplitude}. La \emph{bone integrity} compara el hueso, definido como los vóxeles sobre 150~HU fuera de la geometría metálica, mediante el cambio de volumen y el coeficiente de Sørensen-Dice. La \emph{metal integrity} compara de forma análoga la máscara metálica obtenida de la imagen, con un umbral adaptativo por región. La \emph{streak amplitude} se mide en regiones perpendiculares a las rayas como la diferencia entre el promedio del 5\,\% de desviaciones más altas y el del 5\,\% más bajas respecto de la imagen sin metal del mismo caso. Peters et al. eligen esas regiones de medición sobre las rayas fuertes de las imágenes sin corregir, y declaran como limitación que la elección se hiciera sobre imágenes sin reducción de artefactos, que el propio algoritmo puede alterar \cite{peters2025hybrid}.

Estas métricas se diseñaron para evaluar MAR, no síntesis, y su uso aquí exige invertirlas. La escala de 0 a 4 del protocolo no se traslada, porque se ancla en un algoritmo de MAR (Sección~\ref{sec:obj1}) que no tiene análogo en síntesis. La inversión está definida para la \emph{streak amplitude}, que es el criterio primario, y no para las dos métricas descriptivas \GAPDEC{definición operativa de la inversión de la \emph{bone integrity} y de la \emph{metal integrity} para evaluar síntesis}.

La definición de la \emph{streak amplitude} fija en qué pacientes puede medirse el criterio primario. Su verdad de terreno es la TC original sin metal del mismo paciente, y el campo de desviación es la diferencia en HU entre la imagen sintética y esa original. La partición de prueba reúne 34 pacientes (Sección~\ref{sec:datos}): 20 con implante real y 14 sin metal. La amplitud solo puede medirse en esos 14, porque de un paciente con implante real no hay imagen sin metal. En los 20 restantes lo que se mide es la discrepancia frente a su TC real, con el error absoluto medio en HU dentro de $G$, y se reporta como medida descriptiva.

Las regiones de medición se derivan de la pose y no se trazan a mano. En planos perpendiculares al eje del implante se toman anillos completos, de 360 grados, que van desde una guarda de un vóxel en plano alrededor de $M$ hasta los 12~mm de la banda de generación extendida. Los planos son todos los cortes con máscara del implante menos los 8~mm de cada extremo, la misma exclusión de extremos con que se mide el corredor y el tramo $T$ de SAP (Sección~\ref{sec:sap}). Las regiones quedan así dentro de la banda, que es donde el sintetizador escribe: fuera de ella su amplitud sería nula por construcción. Un estadístico de colas no necesita orientar la región hacia la raya, porque el 5\,\% superior y el 5\,\% inferior encuentran las rayas claras y las oscuras en cualquier punto del anillo. Esta geometría es una desviación declarada del protocolo publicado y no hereda su validación; a cambio, no depende de un lector y es idéntica para los brazos que se comparan.

La guarda se fija en el mínimo por la dirección del sesgo que introduce. Alrededor del metal, el volumen parcial y el artefacto de campo cercano son la misma señal, y separarlos exigiría una imagen sin metal del mismo paciente, que es justo lo que falta donde hay un implante real. Una guarda mayor descartaría señal de artefacto y sesgaría la amplitud a la baja, de modo que se fija en un vóxel en plano, el ancho mínimo que puede ocupar el volumen parcial. Junto a la amplitud se reporta la fracción de vóxeles de la región de medición que quedan exactamente en el suelo de $-1000$~HU de la representación (Sección~\ref{sec:sintetizador}). Si esa fracción es apreciable, la cola inferior queda censurada y la amplitud medida es una cota inferior, con el mismo sentido de sesgo que la guarda.

El realismo se juzga contra las observaciones reales de artefacto (Sección~\ref{sec:datos}), donde la verdad de terreno sin metal no existe. La amplitud de comparación se calcula allí con la forma del estadístico sobre los HU de la banda, y no sobre un campo de desviación, así que no es la misma magnitud que el criterio primario \GAPDEC{qué distancia entre distribuciones de amplitud se usa frente a las observaciones reales de artefacto, y con qué análisis}.

La \emph{streak amplitude} es el único criterio primario, declarado antes de examinar la partición de prueba; las demás métricas son descriptivas. La elección responde a que las rayas son la magnitud que el sintetizador tiene por objeto generar. Fijar un solo criterio primario evita el patrón de comparación múltiple en que basta con que apruebe una de varias métricas. Toda la evaluación del Objetivo~3 es métrica, a diferencia de la del Objetivo~2, que incorpora lecturas humanas del nivel vertebral y del corredor (Secciones~\ref{sec:corredor} y~\ref{sec:marco-metal}) \GAPDEC{si el Objetivo~3 se evalúa solo con métricas, y con qué justificación, o si se añade una lectura humana de la apariencia}.

Sobre ese criterio primario se ordenan tres contrastes, cada uno con lo que puede probar. La comparación de equivalencia frente al protocolo físico es el contraste primario, porque es el único que enfrenta la amplitud sintética con la de un brazo que simula el mecanismo de formación del artefacto. La distancia entre la distribución de amplitudes sintéticas y la medida en las observaciones reales de artefacto es el contraste de realismo. Repite para la apariencia la lógica del Objetivo~2, que compara distribuciones contra una referencia externa y no casos emparejados. La superioridad frente a la inserción por copia y pegado se declara control de cordura y no evidencia de calidad, porque la amplitud de esa línea base vale cero por construcción y cualquier brazo que genere algo la supera. Es una debilidad reconocida del diseño y no una victoria del método. Ningún resultado de ese contraste cuenta a favor del sintetizador, y la decisión del Objetivo~3 recae en la comparación de equivalencia.

Se comparan tres brazos sobre la misma anatomía de origen y las mismas poses de implante: el sintetizador propuesto, la inserción por copia y pegado y el protocolo físico de Peters et al. La inserción por copia y pegado es la línea base ingenua, que coloca metal sin artefacto \GAPDEC{definición operativa de la inserción por copia y pegado, en particular el valor de HU asignado a $M$}. El protocolo físico, implementado con CatSim dentro de XCIST \cite{deman2007catsim,wu2022xcist} y asociado al desafío AAPM CT-MAR \cite{haneda2025aapm}, es el brazo que aporta el mecanismo, porque es el único de los tres que simula la física de formación del artefacto. Se adopta en lugar de una reimplementación validada propia de XCIST/CatSim, que queda fuera de alcance.

El protocolo físico publicado es bidimensional y evalúa MAR, de modo que extenderlo a la síntesis de osteosíntesis pélvica exige una adaptación explícita que no hereda su validación. La geometría con que generan sus datos simula una sola fila de detector con 900 columnas y 1\,000 vistas, con la dispersión calculada para una fila y escalada a 64, y reconstruye por retroproyección filtrada con corrección de agua \cite{peters2025hybrid}. El implante de este trabajo recorre el corredor entero, cuya longitud mediana es de 138~mm (Tabla~\ref{tab:preinscripcion}), así que la adaptación cubre toda la extensión axial del corredor y no un corte aislado. El metal de sus datos de entrenamiento son formas fractales aleatorias insertadas en posiciones aleatorias de tejido blando o hueso, y no implantes \cite{peters2025hybrid}. Peters et al. validan su simulación contra un fantoma físico, sin cubrir el paso híbrido con imágenes clínicas ni esas geometrías genéricas.

Reproducir ese protocolo no es validar el simulador. Validarlo exigiría comparar la simulación contra mediciones físicas en un fantoma con material metálico conocido, escaneado en un tomógrafo real, y esos recursos no están disponibles. Lo que se hizo es una verificación de reproducción: se corrió el protocolo en local y se revisó su código publicado. En ese código, paciente y metal se proyectan juntos, en una sola simulación, sobre un volumen de tres materiales. Son agua con el paciente, agua con la máscara del metal en signo negativo, que desplaza el agua, y la aleación con esa misma máscara. La inserción ocurre, por tanto, por desplazamiento de material antes de la proyección y no por pegado en el dominio de imagen, y eso es lo que hace híbrido al protocolo.

La revisión del código se hizo sobre dos repositorios publicados, en las revisiones \texttt{4cf3544} del simulador y \texttt{4993e87} de los ejemplos, cuya licencia declarada es la de tres cláusulas (BSD 3-Clause). El script de simulación publicado no corre tal como viene: aborta al serializar su propia configuración, de modo que reproducirlo exige un parche de una línea. Ese parche se aplicó sobre una copia local y no sobre el clon del repositorio original, para que el clon revisado quede intacto \GAPLIT{referencia bibliográfica del repositorio publicado del protocolo físico, con las revisiones y la licencia que se revisaron; el código no tiene ficha ni entrada de bibliografía}. El protocolo físico se corre sobre un subconjunto reducido de pacientes \GAPDEC{número de pacientes del subconjunto del protocolo físico}, con un control de simulación y reconstrucción sin metal que aísla los cambios que introduce el propio protocolo \GAPDATO{resultados del protocolo físico de Peters et al. sobre la anatomía y las poses de este trabajo}.

La superioridad sobre la inserción por copia y pegado se evalúa con una prueba de rangos con signo de Wilcoxon de una cola, con nivel de significación de 0.05. La prueba se parea por paciente, porque los dos brazos se aplican a los mismos pacientes. La unidad de agregación es el paciente: la mediana de la amplitud sobre sus regiones de medición y sus cortes deja un valor único. Las pruebas pareadas se calculan sobre los 14 valores de los pacientes de prueba sin metal (Sección~\ref{sec:datos}). Se reportan el tamaño de efecto y su IC además del valor $p$.

Para la comparabilidad con el protocolo físico se usan dos pruebas unilaterales de equivalencia (TOST, \emph{two one-sided tests}). Se concluye equivalencia solo si el IC del 90\,\% de la diferencia pareada cae entero dentro de un margen $[-\Delta, +\Delta]$ fijado de antemano. Nunca se infiere equivalencia de una diferencia no significativa \GAPDEC{si el brazo de equivalencia frente al protocolo físico se mantiene o pasa a trabajo futuro, contingencia de plazo registrada por la autora}. El margen $\Delta$ se define como la variabilidad propia del método bajo condiciones que no deberían cambiar el resultado. Esa variabilidad suma la que producen distintas semillas del sintetizador sobre un mismo caso y la de corridas repetidas del protocolo físico sobre ese caso \GAPDEC{qué cambia entre dos corridas repetidas del protocolo físico}. Ambas se miden solo en los casos de validación de la partición (Sección~\ref{sec:datos}), de modo que el protocolo físico también se corre sobre ellos; tras aplicar el criterio de inclusión, quedan 3 pacientes con implante real. El número de semillas por caso no está fijado (Sección~\ref{sec:sintetizador}).

Fuera de la banda, el RMSE y el índice de similitud estructural (SSIM, \emph{structural similarity index}) contra la TC original comprueban la preservación, restringidos al complemento de $B_{\delta}$. En el sintetizador esa preservación es exacta por construcción (Sección~\ref{sec:sintetizador}); en el protocolo físico sí es una medición, porque la simulación y la reconstrucción pueden alterar la imagen entera. Dentro de $B_{\delta}$, la desviación es la señal buscada y no se puntúa como error.

\subsection{Resumen del diseño experimental}\label{sec:resumen-diseno}

La Tabla~\ref{tab:diseno} liga cada objetivo evaluado con su variable independiente, su variable dependiente, su comparación y su análisis. Los tres fijan de antemano la magnitud que se mide, pero no con el mismo grado de cierre. En el Objetivo~1, el criterio y la regla de decisión se fijaron antes de correr la prueba. En el Objetivo~2 se preinscribieron la distribución de poses y la métrica, pero no una regla de decisión sobre la distancia. En el Objetivo~3 se fijaron el criterio primario, su definición operativa y el orden de los contrastes, y el margen $\Delta$ queda pendiente de preinscripción.

\begin{table}[htbp]
\centering
\small
\caption{Diseño experimental por objetivo.}
\label{tab:diseno}
\begin{tabular}{p{0.12\linewidth} p{0.16\linewidth} p{0.19\linewidth} p{0.17\linewidth} p{0.2\linewidth}}
\hline
\textbf{Objetivo} & \textbf{Variable independiente} & \textbf{Variable dependiente} & \textbf{Comparación} & \textbf{Análisis} \\
\hline
1. Representación & Variante del autoencoder (2) $\times$ codificación (3) & MAE en hueso por paciente, en HU; MAE en $B_{\delta}$, descriptivo & Criterio de 25~HU & Media sobre 34 pacientes de prueba; IC 95\,\% por remuestreo de pacientes; orden a priori \\
2. Colocación & Poses del muestreador preinscrito; cohorte primaria (72) y de sensibilidad (49) & Distribución de grados de brecha (0 a 3); fracción por zona de densidad y viabilidad del corredor por caso, descriptivas & Series navegada y convencional de Zwingmann et al., en S1 & Distancia de Wasserstein-1 en grados, sin regla de decisión; estratificación post hoc por viabilidad del calibre \\
3. Apariencia & Brazo: sintetizador, copia y pegado, protocolo físico & \emph{Streak amplitude} (primaria); \emph{bone} y \emph{metal integrity} (descriptivas) & Protocolo físico; observaciones reales de artefacto; copia y pegado & TOST con margen $\Delta$ de validación, contraste primario; distancia entre distribuciones; Wilcoxon pareado de una cola, control de cordura \\
\hline
\end{tabular}
\end{table}

\section{Amenazas a la validez}\label{sec:amenazas}

Las amenazas se agrupan en cuatro tipos: interna, de constructo, externa y de la conclusión. Para cada una se indica qué se hizo para controlarla o por qué queda abierta.

\textbf{Validez interna.} En el Objetivo~2, la amenaza principal es la circularidad entre muestreador y referencia clínica, y se controla con la preinscripción: ningún parámetro sale de la referencia clínica y la distribución se congeló antes de calcular cualquier distancia contra ella. Tres decisiones, en cambio, se tomaron después de ver datos: la fracción de limpieza de componentes, la regla de ocupación del implante y la estratificación por viabilidad del calibre. El control de nivel excluyó a 11 de 34 pacientes con objetos no ortopédicos y a 7 de 57 sin objeto (Sección~\ref{sec:datos}). La cohorte primaria pierde así, en proporción, más pacientes del grupo~2 que del grupo~3, y su composición por grupos deja de ser la de los 91 pacientes con S1 localizado. Si el corredor difiriera entre grupos, la distribución de grados heredaría esa composición; la cohorte de sensibilidad, que solo reúne pacientes sin objeto, no está expuesta a ese desequilibrio.

La fractura es el confundidor de esa estratificación, y el cribado de la Sección~\ref{sec:datos} lo midió con grupo de comparación. La fractura confirmada aparece por igual en los dos estratos, 10 de 15 corredores estrechos y 10 de 15 controles, así que no explica el estrechamiento. La fractura sacra desplazada, que es la única con mecanismo conocido, sí se concentra en los corredores estrechos: 7 de 15 frente a 2 de 15, con $p = 0.109$. Con 15 casos por grupo esa diferencia es sugerente y no concluyente. El estrechamiento no puede atribuirse entonces ni a anatomía normal ni a patología, y esta muestra no permite cuantificar la fracción que corresponde a cada una.

En el Objetivo~1, la regla operativa y sus seis combinaciones se fijaron después de una exploración del autoencoder preentrenado que incluía a los pacientes de prueba y que ya había fallado el criterio de 25~HU. Ese umbral era anterior a la exploración y no se movió. El rediseño del Objetivo~3 también es posterior al veredicto del Objetivo~1, aunque la regla de la compuerta no se tocó. Lo mismo ocurre con la extensión de la compuerta al latente de Guo et al., cuya cota por saturación se calculó sobre los mismos pacientes de prueba.

\textbf{Validez de constructo.} SAP mide protrusión fuera de una envolvente ósea segmentada, no perforación de una cortical segmentada; su resolución es la del control sobre el fantoma (Sección~\ref{sec:sap}). El cierre morfológico de la envolvente (Sección~\ref{sec:corredor}) podría ocupar el canal sacro o los forámenes, y en ese caso una perforación hacia ellos daría brecha nula \GAPDATO{comprobación de que la envolvente ósea deja fuera el canal sacro y los forámenes sacros}. El marco de referencia se define perpendicular al platillo superior de S1, mientras que el control de nivel y las medidas sobre máscaras usan el techo de la etiqueta de S1 (Sección~\ref{sec:marco-metal}). Ese techo puede estar en la cresta sacra media, así que las dos superficies pueden no coincidir, y esa diferencia no está medida \GAPDATO{diferencia en milímetros entre el techo de la etiqueta de S1, que puede estar en la cresta sacra media, y el platillo superior sobre el que se define el marco de referencia}. Los límites de grado son una convención: el ancho de 2~mm procede de la literatura de tornillos pediculares, de la que Smith et al.~\cite{smith2006iliosacral} declaran tomar la escala, y su traslado al corredor iliosacro es un supuesto explícito.

La referencia clínica no dice a qué grado pertenece una perforación de exactamente 2 o 4~mm, ni reporta el grosor de corte. Tejwani et al.~\cite{tejwani2014} describen más penetración foraminal en cortes de 2.5~mm que en cortes de 5.0~mm, aunque con $p = 0.3$. Cualquier acuerdo con la referencia clínica se lee como acuerdo entre distribuciones de grados medidas a una resolución de referencia desconocida, no como equivalencia vóxel a vóxel. El nivel de la referencia es un supuesto en una de sus dos series. Zwingmann et al. nombran S1 en la técnica de la serie convencional, pero no en la de la navegada. La definición de posición ideal del estudio se refiere al platillo sacro correspondiente y a los forámenes de S1. Este trabajo asume S1 también para la serie navegada, y la publicación no permite verificarlo.

La Ecuación~\eqref{eq:w1} trata los grados como equiespaciados, lo que es también una convención, y la convención $h = 2\sigma$ fija la escala de todas las perturbaciones y, con ella, la distancia resultante. En la apariencia, las métricas de Peters et al. se diseñaron para reducir artefacto y no para generarlo, y una máscara binaria no codifica la aleación del implante. La banda $B_{\delta}$ trunca a propósito las rayas lejanas (Sección~\ref{sec:sintetizador}), y el protocolo físico, que simula y reconstruye la imagen entera, no las trunca. La comparación entre ambos no puede mostrar, por tanto, si el sintetizador reproduciría las rayas más allá de la banda. El sintetizador se entrena, además, con componentes metálicos de todo tipo bajo el supuesto no verificado de que la relación entre máscara y artefacto no depende del tipo de implante (Sección~\ref{sec:sintetizador}). Por último, el trabajo asume que la apariencia de un artefacto metálico puede generarse en el dominio de imagen sin datos de proyección. Lin et al.~\cite{lin2019} y Li et al.~\cite{li2024} reportan una penalización por operar solo en imagen, pero en reducción de artefactos. Lin et al. describen esa reducción como un problema mal planteado, porque paciente y metal contribuyen a la misma traza en el sinograma; la síntesis no tiene que separar esas dos contribuciones. La comparación con el protocolo físico es la que puede contradecir ese supuesto.

El suelo de $-1000$~HU de la representación (Sección~\ref{sec:sintetizador}) censura la cola inferior del campo de desviación, y el sesgo que introduce es siempre a la baja. Ese sesgo no pesa igual en los dos contrastes que miden la amplitud. Frente al protocolo físico apenas pesa, porque su salida reconstruida baja poco por debajo de ese suelo. Frente a la distribución de las observaciones reales de artefacto pesa mucho más, porque esos volúmenes alcanzan valores de TC muy inferiores. Esa diferencia se apoya en una sola corrida local del protocolo físico, de modo que ordena la gravedad del sesgo y no fija una cota. La decisión que trata el suelo compara la pérdida atribuible a la representación con el margen $\Delta$ y lo declara limitación de alcance solo si esa pérdida es menor \GAPDEC{contra cuál de los dos contrastes, el protocolo físico o la distribución de las observaciones reales, se evalúa esa regla}.

\textbf{Validez externa.} La referencia clínica proviene de pelvis fracturadas (tipos~B y~C de Tile), que es la población en la que el procedimiento se practica, mientras que los corredores de este trabajo se miden en pelvis sin osteosíntesis. Reilly et al.~\cite{reilly2003effect} midieron en seis pelvis cadavéricas que un desplazamiento craneal de 5 a 20~mm de una fractura sacra de zona~II, creada por osteotomía, reduce entre 36 y 90\,\% el área disponible para tornillos iliosacros en S1. Una malreducción residual estrecharía, por ese mecanismo, los corredores de la serie clínica. La anatomía receptora tampoco está libre de fractura: el cribado de fractura halló fractura confirmada en 20 de los 30 casos que examinó (Sección~\ref{sec:datos}). Eso acerca la anatomía receptora a la población en que el procedimiento existe, en vez de alejarla. Lo que no puede afirmarse es que las dos poblaciones estén emparejadas, porque la fractura de la cohorte receptora es incidental, no está graduada ni clasificada por un especialista y su estado de reducción es desconocido. Además, la cohorte local es una fracción de la colección publicada, su reconstrucción es desconocida, la publicación de TotalSegmentator no reporta exactitud para la etiqueta de S1, y el protocolo físico adoptado se validó en dos dimensiones y para MAR. En el Objetivo~3, los tres desplazamientos de dominio entre entrenamiento y síntesis (Sección~\ref{sec:sintetizador}) limitan también la transferencia de lo aprendido sobre metal real a los tornillos paramétricos.

\textbf{Validez de la conclusión.} Las dos series de la referencia clínica son pequeñas (26 y 35 tornillos, Sección~\ref{sec:sap}) y cuentan tornillos, no pacientes independientes. La distancia de Wasserstein-1 se reporta sin prueba inferencial, porque la preinscripción no prevé ninguna, y su estabilidad solo se aprecia al comparar las dos cohortes \GAPDEC{si se añade a la distancia de Wasserstein-1 un intervalo por remuestreo de pacientes, como análisis no preinscrito}. En la compuerta del Objetivo~1, la multiplicidad no corregida solo puede favorecer un veredicto aprobatorio, porque basta con que apruebe una de seis combinaciones; como el veredicto fue negativo, esa multiplicidad no lo compromete. La extensión posterior al latente de Guo et al. añade un séptimo candidato, examinado sobre los mismos pacientes de prueba. En la apariencia, el margen $\Delta$ descansa en solo 3 pacientes de validación (Sección~\ref{sec:apariencia}); las semillas multiplican las observaciones por paciente, pero no el número de pacientes. La prueba de equivalencia tiene además muestras pequeñas, acotadas por el subconjunto del protocolo físico (Sección~\ref{sec:apariencia}), de modo que puede no alcanzar a concluir equivalencia ni diferencia.
      % Bloque (c) formal y diseno experimental
\chapter{RESULTADOS Y DISCUSIÓN}
% Bloque (c), parte empirica. Ver redaccion/MAPA.md.
% Pautas de la plantilla UTEC movidas a redaccion/RUBRICA.md (P-RE).

\section{Evaluación de la representación multiventana}

\section{Admisibilidad quirúrgica de la colocación}

\section{Coherencia de apariencia}

\section{Discusión}
      % Bloque (c) empirico: resultados y discusion
\customchapter{CONCLUSIONES} 

Esta sección le brinda la oportunidad al autor de destacar los puntos más importantes de su informe. Para elaborar esta sección puede seguir la siguiente estructura:

\begin{itemize}
  \item \textit{Resumen del aporte:} Re-expresar brevemente el trabajo realizado, los objetivos e hipótesis o preguntas de investigación. Resalte los hallazgos más importantes.
  \item \textit{Evaluación del estudio:} Indique cuáles considera que son los logros y las limitaciones de su trabajo. Evalúe hasta qué punto se han cumplido los objetivos de su investigación.
%  \item \textit{Sugerencias para investigación futura:} Sugiera cómo su trabajo reportado en este documento abre nuevas posibilidades de investigación.
  \item \textit{Implicancias del estudio:} Coloque el estudio en un contexto más amplio de investigación en la disciplina y/o una situación en el mundo real.
  \item \textit{Aplicaciones de la investigación:} Indique cómo la investigación puede ser útil en la práctica en situaciones del mundo real.
\end{itemize}

Un error frecuente es mencionar conclusiones de otros trabajos o autores, por ejemplo, que una técnica $X$ es la más utilizada en la literatura. Concéntrese en discutir sus hallazgos. Pueden usarse viñetas para mostrar las conclusiones más importantes.   % Bloque (d)
\chapter*{\center \Large TRABAJOS FUTUROS} 
\addcontentsline{toc}{section}{\bfseries TRABAJOSFUTUROS} 
\markboth{TRABAJOSFUTUROS}{TRABAJOSFUTUROS} 

Una tesis sin trabajos futuros podría ser interpretado como un trabajo con poco interés por parte del tesista. Sugiera cómo el trabajo reportado en este documento abre nuevas posibilidades de investigación. Nombre algunos puntos que quedaron inconclusos en el proyecto de tesis. También pueden colocarse algunas ideas que, por un problema de tiempos, no formaron parte de este estudio.  % sección opcional

%% ============================================================================
\renewcommand{\bibname}{\hfill\Large\bf{REFERENCIAS BIBLIOGRÁFICAS}\hfill}

\bibliographystyle{IEEEtran} % Estableciendo el estilo de citas IEEEtram.

\bibliography{referencias} % Recibe las referencias de IEEE

\chapter*{\center \Large ANEXOS} 
\addcontentsline{toc}{section}{\bfseries ANEXOS} 
\markboth{ANEXOS}{ANEXOS} 

\par Los algoritmos desarrollados .....

\setcounter{table}{0}
\renewcommand{\thetable}{A.\arabic{table}}

\begin{sidewaystable}
\centering
\footnotesize
\caption{Comparación de los trabajos revisados con este trabajo.}
\label{tab:comparacion}
\begin{tabular}{p{0.10\linewidth} p{0.18\linewidth} p{0.19\linewidth} p{0.18\linewidth} p{0.24\linewidth}}
\hline
\textbf{Trabajo} & \textbf{Enfoque} & \textbf{Colocación del objeto} & \textbf{Fuera del objeto} & \textbf{Resultado reportado y su condición} \\
\hline
Chen et al.~\cite{chen2024tumorsynthesis} & Difusión latente; \emph{inpainting} por máscara de tumor & Máscaras de elipsoides refinadas con radiólogos & No se modela & Dice en hígado de 62.5 a 66.5\,\%, con una de tres redes; cerca de la mitad tomada por real por cuatro radiólogos \\
Ramzan et al.~\cite{ramzan2026claim} & Difusión en el dominio de imagen; \emph{inpainting} & Máscaras muestreadas sobre un modelo clínico de 17 segmentos & Fondo recompuesto con la imagen real en cada paso & Dice de 58.89 a 63.53 en su mejor configuración; fondo evaluado solo de forma cualitativa \\
Jacob et al.~\cite{jacob2026lgesynthnet} & Difusión latente con ControlNet; \emph{inpainting} & Elipsoide en una región elegida al azar & La imagen completa se decodifica del espacio latente & Mejora de Dice al añadir imágenes sintéticas; sus cifras no se citan porque difieren entre el resumen y los resultados del artículo \\
Zhang et al.~\cite{zhang2025diffboost} & ControlNet sobre Stable Diffusion; imagen completa desde ruido & Borde de una máscara de segmentación existente & Cambia la intensidad de toda la imagen, sin efecto ligado al objeto & Dice de 78.46 a 84.56\,\% en próstata, de 94.42 a 94.78\,\% en bazo y de 62.92 a 71.65\,\% en mama \\
Peters et al.~\cite{peters2025hybrid} & Simulación física híbrida, bidimensional & Aleatoria en entrenamiento; sin regla descrita en evaluación & Artefacto simulado en toda la imagen & En el fantoma, valor medio de TC con desviación menor que 2\,\%; discrepancia del ruido de hasta 13.3\,\% en regiones con artefacto, en uno de los experimentos \\
Karageorgos et al.~\cite{karageorgos2024ddpm} & Simulación con CatSim para entrenar MAR por difusión en sinograma & Solape de al menos 50\,\% con una región de más de 200~HU & Artefacto simulado y después eliminado & SSIM de 0.964 en datos simulados; mejor valor en 13 de 28 métricas en cuatro TC clínicas con metal virtual \\
Liu et al.~\cite{liu2025pipeline} & Planificación geométrica por optimización & Una trayectoria por tornillo; sin mención de tornillos iliosacros; no se aplica a fracturas de sacro & No genera imagen & Error de reducción de la fractura de 2.56~mm y $3.31^{\circ}$ en 14 casos; aceptación de 86.7\,\% en 12 casos, según tres cirujanos \\
Zwingmann et al.~\cite{zwingmann2009navigated} & Series clínicas, navegada y convencional & Colocaciones reales: dos distribuciones ordinales & No aplica & Grado~0 de brecha cortical en el 69\,\% (navegada) y el 40\,\% (convencional) de los tornillos \\
Wang et al.~\cite{wang2025adaptiveweighting} & MAR con tres ventanas de HU en cascada & Metal segmentado de datos clínicos e insertado a mano (conjunto dental) & Artefacto eliminado, no generado & En datos simulados y en la ventana estrecha, PSNR de 32.67~dB con supervisión multiventana y de 26.76~dB sin ella; lectura de cinco médicos \\
\hline
Este trabajo & Muestreador de colocación y sintetizador por difusión en el dominio de imagen. Supuesto: el artefacto puede generarse en el dominio de imagen. Convención: ancho de $B_{\delta}$ & Distribución de poses preinscrita alrededor del corredor medido en cada volumen & Genera dentro de $G = M \cup B_{\delta}$ y copia el resto & SAP frente a la referencia clínica, ya ejecutado; sin resultado del sintetizador, que se comparará por \emph{streak amplitude} con la inserción por copia y pegado y, según el plazo, con el protocolo físico \\
\hline
\end{tabular}
\end{sidewaystable}


\end{document}
